\section{Foundations}
\label{ch:foundations}

This chapter fixes the objects that everything else refers to: games, strategies and
values (in the synchronous framework, see Section~\ref{sec:departures}), the
computability conventions, and normal form verifiers. Definitions here are meant to be
final; they are the first formalization targets of the project.

Statements marked \qlib{} are general-purpose quantum information content: they are
intended to be developed in, or contributed to, the
\href{\qliburl}{QuantumLib (Lean-QuantumInfo)} library rather than to remain specific
to this project. The badge marks the general-purpose layer --- games, strategies,
values, and the rigidity theorems of Section~\ref{ch:required-results} --- as opposed
to project-specific machinery such as normal form verifiers and the game
transformations.

\subsection{Games, strategies, and values}

\subsubsection{Games}

\begin{definition}[Nonlocal game]
\label{def:game}
\ledgernode{1.1, 1.1.1}
\lean{MIPRE.Game}
\leanok
\qlib{}
A \emph{two-player one-round game} $\game$ is a tuple
$(\mX, \mY, \mA, \mB, \mu, D)$ where $\mX, \mY$ are finite question alphabets, $\mA,
\mB$ are finite answer alphabets, $\mu$ is a probability distribution on $\mX \times
\mY$, and $D : \mX \times \mY \times \mA \times \mB \to \{0,1\}$ is the decision
predicate ($1$ = accept).
\end{definition}

\begin{definition}[Synchronous game]
\label{def:sync-game}
\uses{def:game}
\lean{MIPRE.SynchronousGame, MIPRE.SynchronousGame.toGame}
\leanok
\qlib{}
A game is \emph{synchronous} if $\mX = \mY$, $\mA = \mB$, and unequal answers to equal
questions are rejected: $D(x,x,a,b) = 0$ whenever $a \neq b$. We write a synchronous
game as $\game = (\mX, \mA, \mu, D)$.
\end{definition}

\subsubsection{Strategies}


\begin{definition}[Tensor-product strategy]
\label{def:tensor-strategy}
\uses{def:game}
\lean{MIPRE.TensorProductStrategy}
\leanok
\qlib{}
A \emph{tensor-product strategy} for $\game$ consists of finite-dimensional Hilbert
spaces $\mH_\alice$, $\mH_\bob$, a unit vector $\ket{\psi} \in \mH_\alice \ot
\mH_\bob$, and projective measurements (PVMs) $\{A^x_a\}$ on $\mH_\alice$ and
$\{B^y_b\}$ on $\mH_\bob$, one for each $x$ (resp.\ $y$): each $A^x_a$ is an
orthogonal projection and $\sum_a A^x_a = \Id$.

Outcome probabilities are computed using the Born rule: the players answer $(a,b)$ to questions $(x,y)$
with probability $\bra{\psi} A^x_a \ot B^y_b \ket{\psi}$. Restricting to projective
measurements follows~\cite{JNVWY20}; by Naimark dilation, allowing general POVMs
would define the same quantum value.
\end{definition}


\begin{definition}[Synchronous strategy]
\label{def:sync-strategy}
\uses{def:sync-game}
\lean{MIPRE.SyncStrategy}
\leanok
\qlib{}
A \emph{synchronous strategy} for a synchronous game $\game = (\mX, \mA, \mu, D)$ is a
dimension $d \geq 1$ together with a family of projective measurements (PVMs)
$\{M^x_a\}_{a \in \mA}$ on $\C^d$, one for each $x \in \mX$: each $M^x_a$ is an
orthogonal projection and $\sum_a M^x_a = \Id$. 

Outcome probabilities are computed with the
tracial state $\tau(M) = \tr(M)/d$: the players answer $(a,b)$ to questions $(x,y)$
with probability $\tau(M^x_a M^y_b)$.
\end{definition}


\begin{definition}[Commuting strategy]
\label{def:commuting-strategy}
\uses{def:sync-game}
\lean{MIPRE.CommutingStrategy}
\leanok
\qlib{}
A \emph{commuting strategy} for a synchronous game $\game = (\mX, \mA, \mu, D)$ is a
unital $C^*$-algebra $\alg$ with a tracial state $\tau$ (a positive linear functional
with $\tau(1) = 1$ and $\tau(ab) = \tau(ba)$ for all $a, b \in \alg$), together with a
family of projective measurements $\{M^x_a\}_{a \in \mA}$ in $\alg$, one for each
$x \in \mX$: each $M^x_a$ is a self-adjoint idempotent and $\sum_a M^x_a = 1$.

Outcome probabilities are computed with $\tau$: the players answer $(a,b)$ to questions
$(x,y)$ with probability $\tau(M^x_a M^y_b)$, a nonnegative real number since
$\tau(M^x_a M^y_b) = \tau(M^x_a M^y_b M^x_a)$.
\end{definition}

\begin{remark}
A synchronous strategy (Definition~\ref{def:sync-strategy}) is exactly the
finite-dimensional case $\alg = \C^{d \times d}$, $\tau = \tr/d$. Via the GNS
representation, commuting strategies capture exactly the synchronous
commuting-operator correlations of the usual two-prover picture, with Bob's
measurements acting in the commutant on $L^2(\alg, \tau)$ --- \cite[Theorem
5.5]{PSSTW16}; see also \cite{Lin23} and \cite[Section 3.4]{Lin25}.
\end{remark}

\begin{definition}[PCC strategy]
\label{def:pcc}
\uses{def:sync-strategy}
\lean{MIPRE.SyncStrategy.IsPCC}
\leanok
\qlib{}
A synchronous strategy $\strategy = (d, M)$ for a synchronous game $\game = (\mX,
\mA, \mu, D)$ is \emph{PCC} (\emph{projective, consistent, and commuting},
following~\cite[Section 5.2]{JNVWY20}) if the measurements associated with any pair
of questions asked with positive probability commute: for all $(x, y)$ in the support
of $\mu$,
\begin{equation*}
  M^x_a M^y_b = M^y_b M^x_a \qquad \text{for all } a, b \in \mA\;.
\end{equation*}
\end{definition}

\begin{remark}[Source PCC and the required representation bridge]
\label{rem:pcc-representation-bridge}
\uses{def:pcc,def:tensor-strategy,lem:sync-le-valstar}
The paper's PCC condition applies to a tensor-product strategy on equal local spaces.
Besides projectivity and $[A^x_a,B^y_b]=0$ on the support of the question distribution,
it requires each measurement separately to be consistent with itself on the state:
\begin{equation*}
  (A^x_a\ot\Id)\ket{\psi}=(\Id\ot A^x_a)\ket{\psi},
  \qquad
  (B^y_b\ot\Id)\ket{\psi}=(\Id\ot B^y_b)\ket{\psi}.
\end{equation*}
These are the conditions of \texttt{def:consistent-measurement},
\texttt{def:consistent-strategy} and \texttt{def:spcc} in the companion
\texttt{games.tex}~\cite{AuditSync26}; see also~\cite[Section 5.2]{JNVWY20}.
The checked realization in Lemma~\ref{lem:sync-le-valstar} uses $A=M$ and
$B=M^{\mathsf T}$. It preserves
probabilities and gives cross-player consistency, but does not establish the displayed
same-measurement identities for arbitrary complex $M$, nor does commutation of two
matrices imply commutation of one with the transpose of the other.

Consequently Definition~\ref{def:pcc} is the current synchronous PCC interface, and
using it as an input to the paper's completeness theorems requires a separate
\emph{PCC representation bridge}. The forward bridge must preserve all probabilities
and support commutation while producing a source PCC strategy. A candidate is to
realify each complex projector on dimension $2d$, then use the same real symmetric
projectors on both halves of a maximally entangled state. Projectivity, the normalized
trace probabilities, the two consistency identities and the symmetry required for
SPCC all need verification for that construction; no such bridge is proved here.
Any transfer back from source PCC strategies to this synchronous interface also needs
its own stated hypotheses and proof. The name PCC in the Lean definition does not
discharge either direction. In both interfaces, commutation is required only on
question pairs in the support of $\mu$ and is distinct from the commuting-operator
model of Definition~\ref{def:commuting-strategy}.
\end{remark}



\subsubsection{Values}

\begin{definition}[Quantum value]
\label{def:q-value}
\uses{def:tensor-strategy}
\lean{MIPRE.TensorProductStrategy.value, MIPRE.quantumValue}
\leanok
\qlib{}
The value of a tensor-product strategy $\strategy = (\mH_\alice,\mH_\bob,\psi,A,B)$ in $\game$ is
\begin{equation*}
  \valstar(\game, \strategy) = \sum_{x, y, a, b} \mu(x,y)\, D(x,y,a,b)\,
  \bra{\psi} A^x_a \ot B^y_b \ket{\psi}\;,
\end{equation*}
and the \emph{quantum value} $\valstar(\game)$ is the supremum of
$\valstar(\game, \strategy)$ over all tensor-product strategies. 
\end{definition}



\begin{definition}[Synchronous value]
\label{def:sync-value}
\uses{def:sync-strategy}
\lean{MIPRE.SyncStrategy.value, MIPRE.syncValue}
\leanok
\qlib{}
The value of a synchronous strategy $\strategy = (d, M)$ in $\game$ is
\begin{equation*}
  \synval(\game, \strategy) = \sum_{x, y, a, b} \mu(x,y)\, D(x,y,a,b)\,
  \tau\bigl(M^x_a M^y_b\bigr)\;,
\end{equation*}
and the \emph{synchronous value} of $\game$ is $\synval(\game) = \sup_\strategy
\synval(\game, \strategy)$.
\end{definition}


\begin{definition}[Commuting value]
\label{def:commuting-value}
\uses{def:commuting-strategy}
\lean{MIPRE.CommutingStrategy.value, MIPRE.commValue}
\leanok
\qlib{}
The value of a commuting strategy $\strategy = (\alg, \tau, M)$ in $\game$ is
\begin{equation*}
  \valco(\game, \strategy) = \sum_{x, y, a, b} \mu(x,y)\, D(x,y,a,b)\,
  \tau\bigl(M^x_a M^y_b\bigr)\;,
\end{equation*}
and the \emph{commuting value} of $\game$ is $\valco(\game) = \sup_\strategy
\valco(\game, \strategy)$.
\end{definition}


\begin{lemma}[A perfect strategy rejects nothing on the support]
\label{lem:perfect-rejects-nothing}
\uses{def:commuting-value,def:sync-value}
\lean{MIPRE.sum_re_tracial, MIPRE.re_nonneg_tracial,
  MIPRE.re_eq_zero_of_tracialValue_eq_one, MIPRE.SyncStrategy.value_eq_tracialValue,
  MIPRE.SyncStrategy.re_eq_zero_of_value_eq_one}
\leanok
Let $\game$ be a synchronous game, $\tau$ a tracial state and $M$ a family of projective
measurements with $\valco(\game, M) = 1$. Then $\tau(M^x_a M^y_b) = 0$ whenever
$\mu(x,y) > 0$ and $D(x,y,a,b) = 0$. In particular a value-$1$ synchronous strategy has
$\mathrm{Tr}(M^x_a M^y_b) = 0$ at every rejected answer pair of every question pair the verifier asks.
\end{lemma}

\begin{proof}
\leanok
At a fixed question pair the outcome probabilities sum to one, $\sum_{a,b} \tau(M^x_a M^y_b)
= \tau(1) = 1$, and each is nonnegative; this is why the value is at most one. The value is
the same sum with the rejected terms deleted, so equality with $1$ forces each deleted term to
vanish, and $\mu(x,y) > 0$ cancels.
\end{proof}

\paragraph{Comments.} Every completeness argument in the pipeline needs this step --- a
strategy of value $1$ does not merely win on average, it assigns probability zero to each
rejected answer at each question actually asked --- and it was used implicitly inside the proof
that the value is at most one, without being available as a lemma. It is stated at the generic
tracial level, so it serves synchronous and commuting strategies alike.

Two formalization notes. The instance set of the ambient algebra must be exactly the
$[\mathrm{Ring}]\,[\mathrm{StarRing}]\,[\mathrm{Module}\ \mathbb{C}]$ of
Definition~\ref{def:commuting-strategy}: demanding a normed algebra instead sends the
elaborator looking for a normed-algebra instance on matrices at every application, which is a
\texttt{whnf} timeout. And the application to a synchronous strategy goes through the named
\texttt{value\_eq\_tracialValue} rather than letting the unifier see through
\texttt{SyncStrategy.value}, for the same reason.

\begin{definition}[Entanglement requirement]
\label{def:ent}
\uses{def:sync-value}
\lean{MIPRE.entRequirement}
\leanok
\qlib{}
For a game $\game$ and $\nu \in [0,1]$, $\Ent(\game, \nu)$ is the least $d$ such that
there exists a synchronous strategy (Definition~\ref{def:sync-strategy}) with dimension at most $d$ achieving value at least
$\nu$ in $\game$, and $\infty$ if no synchronous strategy achieves $\nu$.
\end{definition}


\subsubsection{Relationship lemmas}

\begin{lemma}[Synchronous value lower-bounds the quantum value]
\label{lem:sync-le-valstar}
\uses{def:sync-value,def:q-value}
\lean{MIPRE.syncValue_le_quantumValue}
\leanok
\qlib{}
For every synchronous game $\game$, $\synval(\game) \leq \valstar(\game)$.
\end{lemma}
\begin{proof}
\leanok
Given a synchronous strategy $\{M^x_a\}$ on $\C^d$, take the tensor-product strategy
with $\mH_\alice = \mH_\bob = \C^d$, $A^x_a = M^x_a$, $B^y_b = (M^y_b)^{\mathsf{T}}$,
and $\ket{\psi} = d^{-1/2}\sum_i \ket{i}\ket{i}$ (in Lean,
\texttt{MIPRE.SyncStrategy.toTensorProductStrategy}). On that state
$\bra{\psi} M \ot N \ket{\psi} = \tr(M N^{\mathsf{T}})/d$
(\texttt{MIPRE.dotProduct\_kronecker\_maxEntangled}), so with $N = (M^y_b)^{\mathsf{T}}$ the
Born-rule probability is $\tau(M^x_a M^y_b)$ and the values agree
(\texttt{MIPRE.SyncStrategy.value\_toTensorProductStrategy}). Taking the supremum uses that
the quantum value is a supremum over a possibly empty type with bounded range:
\texttt{MIPRE.TensorProductStrategy.value\_le\_one} and
\texttt{MIPRE.quantumValue\_nonneg}, both proved from the Born rule.
\end{proof}

The converse direction --- every good tensor-product strategy for a synchronous game
whose questions carry proportional weight on the diagonal can be rounded to a good
synchronous strategy --- is the almost-synchronicity theorem, imported as
Theorem~\ref{thm:almost-sync}. The hypothesis on the diagonal cannot be dropped: without
it the two values are unrelated (Remark~\ref{rem:almost-sync-hypothesis}).

\begin{lemma}[Synchronous value lower-bounds the commuting value]
\label{lem:sync-le-valco}
\uses{def:sync-value,def:commuting-value}
\lean{MIPRE.syncValue_le_commValue}
\leanok
\qlib{}
For every synchronous game $\game$, $\synval(\game) \leq \valco(\game)$.
\end{lemma}
\begin{proof}
\leanok
A synchronous strategy $(d, M)$ is a commuting strategy with $\alg = \C^{d \times d}$
and $\tau = \tr/d$, with the same value (in Lean,
\texttt{MIPRE.SyncStrategy.toCommutingStrategy}; the value is preserved by
definitional equality).
\end{proof}

\subsection{Linear constraint system games}
\label{sec:lcs-games}

A well-studied class of nonlocal games arises from systems of linear equations over
$\F_2$: the referee asks Alice for a satisfying assignment to a randomly chosen
equation and Bob for the value of a randomly chosen variable in it, and checks
consistency. The Mermin--Peres magic square is the canonical example. The
formalization of this subsection was contributed by Sean Perazzolo
(\texttt{MIPRE.LCS}); beyond the definitions below it includes observable- and
projector-based strategy formalisms and their equivalence, a sum-of-squares
decomposition of the loss operators, the extraction of operator identities from
perfect play on the EPR state, matrix representations of solution groups, and the
magic square as a worked example
(\texttt{MIPRE.LCS.MagicSquare.merminPeresStrategy}).

\begin{definition}[LCS instance]
\label{def:lcs-instance}
\lean{MIPRE.LCS.Layout, MIPRE.LCS.Game}
\leanok
A \emph{binary linear constraint system} consists of $r$ equations over $s$
$\F_2$-valued variables, specified by the set $V_i \subseteq \{1, \dots, s\}$ of
variables occurring in each equation $i$ together with a right-hand side
$b \in \F_2^r$; equation $i$ reads $\sum_{j \in V_i} x_j = b_i$.
\end{definition}

\begin{definition}[LCS game]
\label{def:lcs-game}
\uses{def:lcs-instance,def:game}
\lean{MIPRE.LCS.Layout.Question, MIPRE.LCS.Layout.Answer,
MIPRE.LCS.Layout.questionDist, MIPRE.LCS.Game.accepts,
MIPRE.LCS.Game.toNonlocalGame}
\leanok
The \emph{nonlocal game} of an LCS instance. Both players receive questions from the same
alphabet --- an equation or a variable --- and answer in the same alphabet --- an assignment
to all the variables, or a single bit. The referee samples an equation $i$ uniformly at
random, a uniformly random variable $j \in V_i$, and a uniform \emph{orientation}: one
player is asked $i$ and the other $j$. The player asked the equation answers an assignment
$a$ (only its restriction to $V_i$ is read), the player asked the variable answers a bit $c$,
and they win if $\sum_{k \in V_i} a_k = b_i$ and $a_j = c$. An answer whose shape does not
match its question loses, as does an off-support pair $j \notin V_i$.
\end{definition}

\paragraph{Comments.} The orientation is what makes the game \emph{symmetric in the players},
and it is not cosmetic. A one-directional version --- the equation always to Alice, the
variable always to Bob --- cannot express a statement about both players' variable
observables, which is exactly what the Pauli basis test consumes from the Magic Square
(Lemma~\ref{lem:ms-direct-anticomm}): Alice would never receive a variable question, so
$A^{\Variable_j}$ would not exist. It also fixes that lemma's constant, the conditional
failures being summed over $2 \sum_i \abs{V_i}$ oriented incidences rather than
$\sum_i \abs{V_i}$ unoriented ones.

The shape and support checks in the decider are the discipline
Definition~\ref{def:lidt-cl} learned the hard way: a missing format check made the seeded CL
game vacuous and a soundness theorem about it false. The formalization pins each outcome with
explicit answers --- an accepted pair in each orientation, and one rejected for disagreement,
one for parity, one for shape, one for being off support.

\begin{definition}[Solution group]
\label{def:solution-group}
\uses{def:lcs-instance}
\lean{MIPRE.LCS.SolutionGroup}
\leanok
The \emph{solution group} of an LCS instance is the group presented by generators
$g_1, \dots, g_s$ and $J$ subject to the relations: all generators are involutions,
$J$ is central, variables occurring in a common equation commute, and
$\prod_{j \in V_i} g_j = J^{b_i}$ for each equation $i$. Operator solutions of the
system correspond to representations sending $J \mapsto -I$ (in Lean,
\texttt{MIPRE.LCS.SolutionGroup.solutionGroupRepresentationOfEPRLoss}).
\end{definition}

\begin{definition}[Binary observables]
\label{def:observable}
\uses{def:lcs-instance}
\lean{MIPRE.LCS.IsObservable, MIPRE.LCS.observableOfMeasurementSystem,
MIPRE.LCS.binary_measurement_eq_projector}
\leanok
An \emph{observable} in a star-algebra is a self-adjoint element squaring to $1$. Binary
projective measurements and observables are two views of the same data: $\{P_0, P_1\}$
gives the observable $P_0 - P_1$, and an observable $O$ gives the measurement
$\{\tfrac12(1 + O), \tfrac12(1 - O)\}$.
\end{definition}

\begin{lemma}[Observables and projective measurements correspond]
\label{lem:observable-projector}
\uses{def:observable}
\lean{MIPRE.LCS.observableToProjector,
MIPRE.LCS.isMeasurementSystem_observableToProjector,
MIPRE.LCS.commute_observableToProjector}
\leanok
For an observable $O$ the family $a \mapsto \tfrac12(1 + (-1)^a O)$ is a projective
measurement indexed by $\F_2$, and the correspondence preserves commutation: if $O_1$ and
$O_2$ commute then so do all of their projectors.
\end{lemma}

\begin{proof}
\leanok
Self-adjointness, idempotence, orthogonality and summation to $1$ are each a two-line
computation from $O^* = O$ and $O^2 = 1$; commutation is inherited term by term.
\end{proof}

\begin{lemma}[From an observable strategy to a projective strategy]
\label{lem:observable-strategy}
\uses{def:observable,lem:observable-projector,def:tensor-strategy}
\lean{MIPRE.LCS.ObservableStrategy.toProjectorStrategy,
MIPRE.LCS.ObservableStrategy.isMeasurementSystem_aliceMeasurement,
MIPRE.LCS.ObservableStrategy.aliceMeasurement_commute_bobMeasurement}
\leanok
A strategy given by commuting observables, one per variable of an LCS instance, yields a
strategy given by genuine projective measurements: for a question consisting of several
variables the joint measurement is the product of the corresponding binary projectors,
which is projective because the observables commute, and the two players' measurements
commute with each other.
\end{lemma}

\begin{proof}
\leanok
The joint family over an equation's support is a product of commuting binary projectors,
so it is a projective measurement and sums to $1$ over the outcomes; commutation across
players is inherited from commutation of the underlying observables.
\end{proof}

\begin{definition}[The Magic Square]
\label{def:magic-square}
\uses{def:lcs-instance,def:lcs-game,def:solution-group}
\lean{MIPRE.LCS.MagicSquare.layout, MIPRE.LCS.MagicSquare.game,
MIPRE.LCS.MagicSquare.solutionGroup}
\leanok
The Magic Square is the LCS instance over $\F_2$ with nine variables in a $3 \times 3$
grid and six equations, one per row and column, the rows summing to $0$ and the columns to
$1$. It has no classical solution, since the six equations sum to $0 = 1$.
\end{definition}

\begin{lemma}[The Mermin--Peres strategy is perfect]
\label{lem:mermin-peres}
\uses{def:magic-square,def:observable,lem:observable-strategy}
\lean{MIPRE.LCS.MagicSquare.grid, MIPRE.LCS.MagicSquare.isObservable_grid,
MIPRE.LCS.MagicSquare.grid_sameEquation_comm, MIPRE.LCS.MagicSquare.merminPeresStrategy}
\leanok
There are nine observables on $(\C^2)^{\ot 2}$, one per cell of the Magic Square, such
that the three in each row or column commute and multiply to $+1$ for a row and $-1$ for a
column. They give a strategy for the Magic Square game of value $1$ on two EPR pairs.
\end{lemma}

\begin{proof}
\leanok
The grid is the Mermin--Peres square of tensor products of Pauli matrices. That each entry
is an observable, and that entries sharing an equation commute, are finite computations on
$4 \times 4$ matrices.
\end{proof}

\begin{theorem}[Perfect strategies from operator solutions]
\label{thm:lcs-perfect}
\uses{def:lcs-game,def:tensor-strategy,def:q-value}
\lean{MIPRE.LCS.exists_tensorStrategy_value_eq_one_of_localLoss_annihilates_epr}
\leanok
If a bipartite observable strategy for an LCS instance is perfect --- its local loss
operators annihilate the EPR vector --- then the associated LCS game has a
tensor-product strategy of value $1$.
\end{theorem}
\begin{proof}
Take both players' spaces to be $\C^n$ with the normalized EPR state; Alice measures
with the joint projective measurements of her equations and Bob with the binary
measurements of his observables. The Born-rule value then equals the EPR expectation
of the winning operator, which is $1$ by hypothesis.
\end{proof}

\subsection{Distance measures}
\label{sec:distances}

\begin{definition}[Normalized Hilbert--Schmidt norm]
\label{def:hs-norm}
\lean{MIPRE.hsInner, MIPRE.hsNormSq}
\leanok
\qlib{}
For matrices $A, B \in \C^{d \times d}$, let
$\langle A, B \rangle_{hs} = \tau(A^\dagger B)$ and
$\norm{A}_{hs}^2 = \tau(A^\dagger A)$, where $\tau = \tr/d$ is the
dimension-normalized trace.
\end{definition}

\begin{definition}[POVM distance]
\label{def:distance}
\ledgernode{1.1.1.1}
\uses{def:game,def:hs-norm}
\lean{MIPRE.povmDistance, MIPRE.IsPOVMClose}
\leanok
\qlib{}
Let $\mu$ be a distribution on a finite set $\mX$, and
$\{M^x_a\}$, $\{N^x_a\}$ families of POVMs. We say $M$ and $N$ are
\emph{$\delta$-close}, written $M^x_a \approx_\delta N^x_a$, if
$\E_{x \sim \mu} \sum_a \norm{M^x_a - N^x_a}_{hs}^2 \leq \delta$. Since each
difference $M^x_a - N^x_a$ is Hermitian, the summand equals
$\tau((M^x_a - N^x_a)^2)$.
\end{definition}

\begin{definition}[State-dependent POVM distance]
\label{def:state-distance}
\uses{def:game}
\lean{MIPRE.stateVec, MIPRE.stateNorm, MIPRE.stateSqNorm, MIPRE.stateDist,
  MIPRE.IsStateClose, MIPRE.povmStateDist, MIPRE.IsPOVMStateClose,
  MIPRE.quadForm_eq, MIPRE.stateSqNorm_eq}
\leanok
Let $\mu$ be a distribution on a finite set $\mX$, let
$\ket{\psi} \in \C^{d_\alice} \ot \C^{d_\bob}$, and let $\{A^x\}$, $\{B^x\}$ be families of
operators on $\C^{d_\alice}$. We say $A$ and $B$ are \emph{$\delta$-close on $\ket{\psi}$},
written $A^x \approx_\delta B^x$ when $\mu$ and $\ket{\psi}$ are clear from context, if
\begin{equation*}
  \E_{x \sim \mu} \bra{\psi} \bigl( (A^x - B^x) \ot \Id_\bob \bigr)^\dagger
    \bigl( (A^x - B^x) \ot \Id_\bob \bigr) \ket{\psi} \;\leq\; \delta\;,
\end{equation*}
equivalently if $\E_{x\sim\mu} \norm{((A^x - B^x) \ot \Id_\bob)\ket{\psi}}^2 \leq \delta$. For
families indexed by answers as well, $\{A^x_a\}$ and $\{B^x_a\}$, the outcome sum sits inside the
question average.
\end{definition}

\paragraph{Comments.} This is \emph{not} Definition~\ref{def:distance}, and the difference is
not presentational. Definition~\ref{def:distance} is the normalized Hilbert--Schmidt distance,
which is the distance of a \emph{synchronous} strategy --- there the state is the normalized
trace itself, which is why \texttt{MIPRE/Foundations/Closeness.lean} is the calculus of that one.
The appendix of the Pauli basis test (Section~\ref{sec:rr-qld}) works throughout with the
state-dependent distance above, on a bipartite state that is fixed by context, and every
$\approx_\delta$ in that section is this one.

Two departures from the paper's own definition (\cite{AuditSync26}, \texttt{qld-prelim.tex}),
both in the direction of a stronger statement. The paper writes $\leq O(\delta)$ \emph{inside}
the relation; here the relation is an inequality against $\delta$ and the constants appear in
the lemmas that produce them, so Lemma~\ref{lem:qld-averaging} preserves $\delta$ exactly and
Lemma~\ref{lem:qld-povm-to-obs} gives exactly $\abs{\mA}\delta$. And the paper defines
$\approx_\delta$ only for question-indexed operators, using the answer-indexed form without
comment; both are defined here.

\begin{definition}[Cross-party state-dependent distance]
\label{def:cross-distance}
\uses{def:game,def:state-distance}
\lean{MIPRE.xSqNorm, MIPRE.xPovmDist, MIPRE.IsXPOVMClose, MIPRE.POVM.mul_self_le_self,
  MIPRE.sum_stateSqNorm_le_one, MIPRE.sum_stateSqNormB_le_one, MIPRE.sub_kronecker_right,
  MIPRE.xSqNorm_sum_le_two_mul, MIPRE.sum_bornProb_mapA, MIPRE.sum_bornProb_mapB,
  MIPRE.sum_weight_bornProb_map, MIPRE.sum_bornProb_map,
  MIPRE.one_sub_sum_bornProb_le_condFail, MIPRE.xSqNorm_sum_le_condFail,
  MIPRE.sum_mul_condFail_le, MIPRE.sum_condFail_le_of_pushforward,
  MIPRE.sum_condFail_le_of_le, MIPRE.xStateDist, MIPRE.xStateDist_obsOf_le,
  MIPRE.stateVecB_smul, MIPRE.stateVecB_sum, MIPRE.xNorm, MIPRE.xNorm_nonneg,
  MIPRE.xSqNorm_eq_sq, MIPRE.xNorm_eq_snorm, MIPRE.norm_stateVec_eq_snorm,
  MIPRE.norm_stateVecB_eq_snorm, MIPRE.xNorm_mul_le, MIPRE.xNorm_mul_le_one,
  MIPRE.xSqNorm_mul_le_one, MIPRE.norm_stateVec_anticomm_le, MIPRE.obs2,
  MIPRE.obsOf_sgn_eq_obs2, MIPRE.obs2_conjTranspose, MIPRE.obs2_mul_self_le_one,
  MIPRE.norm_stateVec_comm_le, MIPRE.xSqNorm_eq_snorm_sq, MIPRE.sum_weighted_xSqNorm_eq,
  MIPRE.sum_weighted_xSqNorm_eq', MIPRE.stateSqNorm_sum_le, MIPRE.stateSqNorm_sub_comm,
  MIPRE.norm_stateVecB_sub_comm, MIPRE.mulVec_kronecker_swapVec, MIPRE.swapVec_sub,
  MIPRE.norm_swapVec, MIPRE.snorm_swapVec_aOp_sub_bOp, MIPRE.snorm_swapVec_aOp}
\leanok
Let $\mu$ be a distribution on a finite set $\mX$, let
$\ket{\psi} \in \C^{d_\alice} \ot \C^{d_\bob}$, and let $\{M^x_a\}$ and $\{N^x_a\}$ be families
of POVMs with a common question set $\mX$ and a common answer set, the first on
$\C^{d_\alice}$ and the second on $\C^{d_\bob}$. We say they are \emph{$\delta$-close across the
two parties}, written $M^x_a \ot \Id \simeq_\delta \Id \ot N^x_a$, if
\begin{equation*}
  \E_{x \sim \mu} \sum_a \norm{\bigl( M^x_a \ot \Id_\bob - \Id_\alice \ot N^x_a \bigr)
    \ket{\psi}}^2 \;\leq\; \delta\;.
\end{equation*}
Then, for any two POVMs $M$ and $N$ with a common answer set and $\ket{\psi}$ a unit vector,
\begin{equation}
  \label{eq:cross-from-agreement}
  \sum_a \norm{\bigl( M_a \ot \Id - \Id \ot N_a \bigr)\ket{\psi}}^2
    \;\leq\; 2 \Bigl( 1 - \sum_a \bra{\psi} M_a \ot N_a \ket{\psi} \Bigr)\;,
\end{equation}
twice the probability that the two measurements disagree.
\end{definition}

\paragraph{Comments.} This is the $\simeq_\delta$ of the Pauli appendix, and it is not
Definition~\ref{def:state-distance}: that one compares two families of operators on the
\emph{same} factor, which is what the orthonormalization step of
Corollary~\ref{cor:ortho-from-consistency} needs, whereas every consistency statement of
Lemma~\ref{lem:qld-win} compares one player's operator with the \emph{other} player's. The paper
uses both symbols and says the two are interchangeable up to
Fact~\texttt{fact:agreement}; the formalization keeps them apart, because the conversion is the
factor $2$ of~\eqref{eq:cross-from-agreement} and it is better spent once than assumed
throughout.

Inequality~\eqref{eq:cross-from-agreement} is the whole analytic content of the file
(\texttt{MIPRE/Foundations/CrossConsistency.lean}): expand each square, and the two diagonal
sums are at most one because a POVM element satisfies $A^\dagger A = A^2 \leq A$
(\texttt{POVM.mul\_self\_le\_self}) and the elements sum to the identity. \textbf{No projectivity
is used}, which matters --- the strategies it is applied to are arbitrary POVM strategies, and
projectivity arrives only later, through Corollary~\ref{cor:ortho-from-consistency}.

Two things the formalization needs that the paper leaves implicit. First, what a subtest buys is
not that its decider \emph{equals} an agreement predicate but that \emph{accepting implies
agreeing}: a real decider checks the answer formats first and rejects an ill-formatted pair even
when the post-processings agree. So the readings that coarse-grain the answers may send an
ill-formatted answer anywhere --- $0$, here --- and the implication is all
\texttt{one\_sub\_sum\_bornProb\_le\_condFail} asks. Second, ``divided by the probability that
the subtest is selected'' needs no injectivity: the subtests of the Pauli basis test are indexed
by the verifier's content and several contents give the same question pair, a $(\pauli, W)$
question reading none of it, so the hypothesis is a condition on the push-forward
(\texttt{sum\_condFail\_le\_of\_pushforward}).

\texttt{xStateDist\_obsOf\_le} is Lemma~\ref{lem:qld-povm-to-obs} across the two parties, with
the same proof and the same factor $\abs{\mA}$: the appendix needs it to pass from a two-outcome
coarse-graining of a point measurement to the $\pm 1$-observable it defines, and
Definition~\ref{def:state-distance}'s version compares the wrong pair.

\textbf{The order reversal, and where the sign comes from.} One algebraic rule is proved here
rather than inlined at its two uses, because the paper's own correction note on
Lemma~\ref{lem:qld-obs-commutation} records two displays that had it backwards. If
$A \ot \Id \simeq \Id \ot A'$ and $B \ot \Id \simeq \Id \ot B'$ then
\begin{equation}
  \label{eq:order-reversal-bp}
  AB \ot \Id \;\simeq\; \Id \ot B'A'\;,
\end{equation}
Bob's factors in the \emph{reverse} order, at the sum of the two deviations
(\texttt{xNorm\_mul\_le\_one}; \texttt{xNorm\_mul\_le} carries operator-norm bounds $K$ and $L$ in
front of them). The proof is the split
$AB \ot \Id - \Id \ot B'A' = (A \ot \Id)(B \ot \Id - \Id \ot B') + (\Id \ot B')(A \ot \Id - \Id
\ot A')$, and the reversal is exactly the step that commutes $A \ot \Id$ past $\Id \ot B'$ to
factor the second summand: it puts $B'$ to the \emph{left} of $A'$. Applying it in both orders
and cancelling Bob's two products gives the form the anticommuting case of
Lemma~\ref{lem:qld-obs-commutation} consumes,
\begin{equation}
  \label{eq:anticomm-transfer-bp}
  \norm{(AB + BA) \ot \Id \ket{\psi}} \;\leq\; 2 \norm{(A \ot \Id - \Id \ot A')\ket{\psi}}
    + 2 \norm{(B \ot \Id - \Id \ot B')\ket{\psi}}
    + \norm{\Id \ot (A'B' + B'A') \ket{\psi}}\;,
\end{equation}
\texttt{norm\_stateVec\_anticomm\_le}: an anticommutation on Bob's side transfers to Alice's. The
two reversals do not cancel each other, because what sits between them is a \emph{sum}. Both
statements are on the unsquared deviation, since the triangle inequality is what the argument
uses; \texttt{xSqNorm\_mul\_le\_one} squares the first at the usual cost of a factor two.

Each operator in sight has to be a contraction, and for the $\pm 1$-observable of a two-outcome
POVM that is \texttt{POVM.sub\_mul\_self\_le\_one}, which needs the observable written as a
\emph{difference} of two POVM elements rather than as the weighted sum $\sum_a (-1)^a M_a$ that
\texttt{obsOf} produces. \texttt{obs2} is that difference and \texttt{obsOf\_sgn\_eq\_obs2}
identifies the two, so a consumer can quote whichever form its input is stated in.

\textbf{The swap, and the commutator form.} Two more entries serve the appendix's commutation
step. `norm_stateVec_comm_le` is~\eqref{eq:anticomm-transfer-bp} with a commutator in the middle
instead of an anticommutator --- the same split, Bob's two products cancelling either way. And
`snorm_swapVec_aOp_sub_bOp` exchanges the two sides of a cross-party deviation, which is what lets
a lemma proved for Alice's factor be applied to Bob's: the clean statement is at the level of
vectors, $(X \ot Y)\,\mathrm{swap}\ket{\psi} = \mathrm{swap}\,(Y \ot X)\ket{\psi}$, and everything
else follows by linearity. Corollary~\ref{cor:ortho-from-consistency} produces a projective
measurement on \emph{Alice's} factor only, so the commutation analysis below is run on Bob's
observables and carried back; that swap is the bridge.

\subsection{The expanded state}

\begin{definition}[The expanded state, and an inert ancilla]
\label{def:expanded-state}
\uses{def:cross-distance,def:state-distance}
\lean{MIPRE.expVec, MIPRE.expEquiv, MIPRE.sum_prod_mul, MIPRE.norm_evec_expVec,
  MIPRE.expVec_dotProduct, MIPRE.expVec_unit, MIPRE.mulVec_kron_expVec,
  MIPRE.norm_stateVec_kron_unitary, MIPRE.norm_stateVec_expVec_kron_one,
  MIPRE.expVec_dotProduct_of_pair, MIPRE.mulVec_kron_kron_expVec, MIPRE.bornProb_expVec_kron,
  MIPRE.POVM.kron, MIPRE.POVM.kron_mats, MIPRE.bornProb_sum_sum, MIPRE.sum_bornProb_le_map,
  MIPRE.IsPVM.toPOVM, MIPRE.IsPVM.toPOVM_mats, MIPRE.POVM.ext', MIPRE.POVM.map_mats,
  MIPRE.POVM.map_map, MIPRE.POVM.kron_map, MIPRE.IsPVM.coarse, MIPRE.isPVM_kron,
  MIPRE.isPVM_povm_map, MIPRE.isPVM_povm_kron}
\leanok
Let $\ket{\psi}$ be a state on $\C^{d_\alice} \ot \C^{d_\bob}$ and $\ket{e}$ one on
$\C^{a} \ot \C^{a'}$. The \emph{expanded state} is $\ket{\psi} \ot \ket{e}$, \textbf{re-bipartitioned}
so that the first party holds $\C^{d_\alice} \ot \C^{a}$ and the second $\C^{d_\bob} \ot \C^{a'}$:
one half of the ancilla goes to each. Then
\begin{enumerate}
\item an operator of product form acts factor by factor,
  $\bigl((M \ot N) \ot \Id\bigr)(\ket{\psi} \ot \ket{e})
   = \bigl((M \ot \Id)\ket{\psi}\bigr) \ot \bigl((N \ot \Id)\ket{e}\bigr)$;
\item the norm of a product state is the product of the norms, so the expanded state is a unit
  vector when both factors are;
\item consequently an operator with an \emph{inert} ancilla has the same state-norm on the expanded
  state as on the original one: $\norm{(M \ot \Id_a) \ot \Id \,\ket{\hat\psi}} =
  \norm{(M \ot \Id_\bob)\ket{\psi}}$ when $\ket{e}$ is a unit vector;
\item and a \emph{unitary} on the ancilla is invisible:
  $\norm{(M \ot U) \ot \Id \ket{\hat\psi}} = \norm{(M \ot \Id_a) \ot \Id \ket{\hat\psi}}$.
\end{enumerate}
\end{definition}

\paragraph{Comments.} Formalized in \texttt{MIPRE/Foundations/Expanded.lean}. This is the
bookkeeping of the appendix's \texttt{sec:expanding}, isolated from the Pauli test: items 3 and 4 are
what let a bound proved \emph{before} the expansion be used \emph{after} it, and they are why the
expansion stage costs exactly the constant of Lemma~\ref{lem:qld-obs-commutation} and nothing more.
The re-bipartition is the point --- the ancilla halves are split between the parties, which is what
makes the ancilla observables tied together --- and it is why the statements are about
$\ket{\psi} \ot \ket{e}$ rather than about a tensor power.

\textbf{Born probabilities factorize, and that is where data processing lives.} The consistency
half of the expansion stage is about a measurement of \emph{product} form, and what it needs is
that its Born probabilities split: $\bra{\hat\psi} (M \ot N) \ot (M' \ot N') \ket{\hat\psi} =
\bra{\psi} M \ot M' \ket{\psi} \cdot \bra{e} N \ot N' \ket{e}$
(\texttt{bornProb\_expVec\_kron}; only the ancilla's two operators need be positive, which is what
makes their quadratic form real so the real part of the product splits). \texttt{POVM.kron} is the
product measurement, and \texttt{sum\_bornProb\_le\_map} is the data-processing inequality that
holds: coarse-graining \emph{both} players the same way can only increase the probability that they
agree.

That last one is stated at the Born level deliberately. The state-dependent distance has
\textbf{no} data-processing inequality --- NW19 states theirs for $\simeq_\delta$ and the remark
following it gives a counterexample for $\approx_\delta$ --- so every coarse-graining in the
appendix is taken on Born probabilities and converted once, at the end, by
\eqref{eq:cross-from-agreement}. A formalization that reached for the $\approx$ form would be
proving something false.

\subsection{Naimark dilation against a fixed state}

\begin{lemma}[Naimark dilation of a question-indexed family]
\label{lem:naimark-dilation}
\uses{def:game}
\lean{MIPRE.ancillaProj, MIPRE.ancillaProj_conjTranspose, MIPRE.ancillaProj_mul_self,
  MIPRE.sum_ancillaProj, MIPRE.exists_isometry_of_povm, MIPRE.ancillaEmbed,
  MIPRE.ancillaEmbed_isometry, MIPRE.exists_unitary_extending,
  MIPRE.exists_projective_dilation, MIPRE.dotProduct_mulVec_submatrix,
  MIPRE.dotProduct_comp_equiv, MIPRE.dotProduct_mulVec_conj}
\leanok
Let $\{E^x_a\}$ be a family of POVMs on $\C^d$, indexed by questions $x$ and outcomes
$a \in \mA$. Then there is a family of \emph{projective} measurements $\{P^x_a\}$ on
$\C^d \ot \C^{\mA}$ and a \emph{single}, question-independent isometry
$V : \C^d \to \C^d \ot \C^{\mA}$, $\ket{v} \mapsto \ket{v} \ot \ket{a_0}$, such that
$V^\dagger P^x_a V = E^x_a$ for every $x$ and $a$.
\end{lemma}

\begin{proof}
\leanok
The columns of the matrix $V^x$ whose $(q, j)$ entry is $\sqrt{E^x}$-like --- concretely, the
square root of $E^x_a$ read off in the $a$-th block --- form an isometry $\C^d \to \C^d \ot
\C^{\mA}$ compressing the block-diagonal projectors $\Pi_a = \Id \ot \ket{a}\bra{a}$
back to $E^x_a$. Extending each $V^x$ to a unitary $U^x$ of the larger space that agrees with it on
the $a_0$-slice --- the completion of an orthonormal family to an orthonormal basis --- makes
$P^x_a = (U^x)^\dagger \Pi_a U^x$ projective with the compression along the
\emph{fixed} isometry $\ket{v} \mapsto \ket{v} \ot \ket{a_0}$. The question-independence of that
isometry is the whole point: the expectations are taken against a shared state, which must not
depend on the question.
\end{proof}

\paragraph{Comments.} Formalized in \texttt{MIPRE/Foundations/Dilation.lean}, extracted from the
repetition development where it was written for \texttt{lem:povm-value-eq}. Its other consumers are
Lemma~\ref{lem:ms-direct-anticomm}, which dilates the Magic Square's six constraint POVMs against
one shared ancilla, and Theorem~\ref{thm:linearity}.

\subsection{The commutation analysis}

\begin{lemma}[Commutation analysis]
\label{lem:commutation-analysis}
\uses{def:cross-distance,def:game}
\lean{MIPRE.qform_nonneg_of_nonneg, MIPRE.qform_le_of_le, MIPRE.sum_snorm_sq_mul_le,
  MIPRE.sum_snorm_sq_triangle', MIPRE.sum_snorm_sq_triangle3,
  MIPRE.sum_weighted_snorm_sq_triangle3, MIPRE.POVM.sum_mats_map_prod,
  MIPRE.POVM.sum_mats_map_prod', MIPRE.sum_xSqNorm_le_of_two_step, MIPRE.aOp_mono,
  MIPRE.bOp_mono, MIPRE.bOp_sum, MIPRE.sum_aOp_conjTranspose_mul_self_le_one,
  MIPRE.sum_bOp_conjTranspose_mul_self_le_one, MIPRE.sum_bOp_conjTranspose_mul_self_of_isPVM,
  MIPRE.IsPVM.marg_mul_marg, MIPRE.IsPVM.marg_mul_marg', MIPRE.IsPVM.sum_marg_left,
  MIPRE.IsPVM.sum_marg_right, MIPRE.IsPVM.marg_left, MIPRE.IsPVM.marg_right,
  MIPRE.sum_prod_snorm_sq_mul_le_fst, MIPRE.sum_prod_snorm_sq_mul_le_snd,
  MIPRE.commutation_analysis_abstract, MIPRE.commutation_analysis,
  MIPRE.commutation_analysis_aOp, MIPRE.obs2_commutator_eq}
\leanok
Let $\{A_b\}$ and $\{C_c\}$ be POVMs on $\C^{d_\alice}$ and let $\{B_{b,c}\}$ be a
\emph{projective} measurement on $\C^{d_\bob}$ with outcomes the pairs, with marginals
$B_b = \sum_c B_{b,c}$ and $B_c = \sum_b B_{b,c}$. If
\begin{equation*}
  A_b \ot \Id \simeq_\delta \Id \ot B_b \qquad\text{and}\qquad
  C_c \ot \Id \simeq_\delta \Id \ot B_c\;,
\end{equation*}
then Alice's two POVMs approximately commute on the state:
\begin{equation}
  \label{eq:commutation-analysis-bp}
  \sum_{b, c} \norm{\bigl( [A_b, C_c] \ot \Id \bigr) \ket{\psi}}^2 \;\leq\; 16\, \delta\;.
\end{equation}
\end{lemma}

\begin{proof}
\leanok
\uses{def:cross-distance}
Both products reach the \emph{same} operator on Bob's factor, $\Id \ot B_{b,c}$. For
$A_b C_c \ot \Id$: put $A_b \ot \Id$ in front of $C$'s relation, and put $\Id \ot B_c$ in front of
$A$'s; the two meet at $A_b \ot B_c$, and the second step lands on
$\Id \ot B_c B_b = \Id \ot B_{b,c}$ --- which is where projectivity is used, and the only place
(\texttt{IsPVM.marg\_mul\_marg}). For $C_c A_b \ot \Id$, the same two steps in the other order,
using $B_b B_c = B_{b,c}$. Each side is therefore within $4\delta$ after one triangle inequality,
and~\eqref{eq:commutation-analysis-bp} is one more.

Putting an operator in front costs nothing and \emph{adds} its index to the sum, which is the
paper's \texttt{fact:add-a-proj}: if $\sum_i F_i^\dagger F_i \leq \Id$ then
$\sum_i \norm{F_i v}^2 \leq \norm{v}^2$. For a POVM that hypothesis is $A_b^2 \leq A_b$ with
$\sum_b A_b = \Id$, so it needs no projectivity of Alice's measurements
(\texttt{sum\_aOp\_conjTranspose\_mul\_self\_le\_one}).
\end{proof}

\paragraph{Comments.} Formalized in \texttt{MIPRE/Foundations/Commutation.lean}. This is the
paper's \texttt{lem:commutation-analysis} (\texttt{games.tex}), and it is the only step of the
Pauli appendix that genuinely uses projectivity of a measurement.

\textbf{Proved for abstract families in one algebra.} The two parties enter only through four
hypotheses --- the two factors commute, and Bob's two marginals multiply to the joint element in
either order --- so \texttt{commutation\_analysis\_abstract} is stated for families in a single
matrix algebra and the bipartite version is its instance. That is not tidiness: stating it with
$\ot$ in the hypotheses made the unifier's work quadratic in the number of rewrites and the proof
timed out; in one algebra it elaborates in seconds.

\textbf{No square root, and the constant is exact.} Every step is a triangle inequality or an
operator-norm bound, so $16\delta$ is what the argument gives. The last step of the appendix's
commuting case is also exact: expanding a two-outcome observable as $\Id - 2 M_1$, the commutator
of two of them is \emph{four times} the commutator of the two $M_1$'s
(\texttt{obs2\_commutator\_eq}), because the identity commutes with everything --- a factor $16$ in
the squared norm and no approximation.

\subsection{Attaching a projection, and the pasting lemma}

\begin{lemma}[Cool closeness fact, and the marginal step]
\label{lem:cool-closeness-fact}
\uses{def:state-distance,def:cross-distance}
\lean{MIPRE.snorm_sq_mul_le_of_contraction, MIPRE.snorm_sq_sum_proj_mul,
  MIPRE.sum_snorm_sq_cool, MIPRE.sum_fibre_fst, MIPRE.sum_snorm_sq_cool_prod,
  MIPRE.IsPVM.aOp, MIPRE.IsPVM.bOp, MIPRE.sum_xSqNorm_marg_le,
  MIPRE.sum_snorm_sq_mul_proj_le, MIPRE.sum_xSqNorm_marg_le'}
\leanok
Let $\{A_i\}_{i \in \mA}$ be a projective measurement and $\{B_i\}$ any family of
operators on the same space, with
$\sum_i \norm{(A_i - B_i)\ket{\psi}}^2 \leq \delta$. Then for any map $f : \mA \to \mB$,
\begin{equation}
  \label{eq:cool-closeness-bp}
  \sum_{k} \Bigl\| \sum_{i \,:\, f(i) = k} \bigl( A_i - A_i B_i \bigr) \ket{\psi} \Bigr\|^2
    \;\leq\; \delta\;,
\end{equation}
\emph{simultaneously} over all the fibres of $f$ --- multiplying each element by $A$'s own element
costs nothing, and neither does summing along a fibre.

Consequently, if $\{Q_{a,b}\}$ is a projective measurement on $\C^{d_\alice}$ with outcomes pairs,
$\{Z_b\}$ is a projective measurement on $\C^{d_\bob}$, and $\{X_a\}$ is any family on
$\C^{d_\bob}$, then the $b$-marginal of $Q$ inherits $Q$'s consistency with the ordered product:
\begin{equation}
  \label{eq:marginal-step-bp}
  \sum_a \norm{\Bigl( \bigl(\sum_b Q_{a,b}\bigr) \ot \Id - \Id \ot X_a \Bigr)
      \ket{\psi}}^2
    \;\leq\; 2 \sum_{a, b} \norm{\bigl( Q_{a,b} \ot Z_b - Q_{a,b} \ot \Id \bigr)\ket{\psi}}^2
      + 2 \sum_{a, b} \norm{\bigl( Q_{a,b} \ot \Id - \Id \ot Z_b X_a \bigr)\ket{\psi}}^2\;.
\end{equation}
\end{lemma}

\begin{proof}
\leanok
\uses{def:state-distance}
For~\eqref{eq:cool-closeness-bp}: projectivity gives $A_i - A_i B_i = A_i (A_i - B_i)$, and the
cross terms of one fibre vanish because $A_i A_{i'} = 0$ for $i \neq i'$, so the fibre's squared
norm is the sum of its terms' squared norms. Each term loses only $A_i \leq \Id$, and the fibres
are disjoint, so the outer sum is absorbed rather than repeated.

For~\eqref{eq:marginal-step-bp}, the intermediate operator is
$\sum_b Q_{a,b} \ot Z_b$. Reaching it from $\sum_b Q_{a,b} \ot \Id$ is exactly the first part,
applied along the fibres of the first projection. Reaching $\Id \ot X_a$ from it uses
$\sum_b Z_b = \Id$ to write the difference as
$\sum_b (\Id \ot Z_b)\bigl( Q_{a,b} \ot \Id - \Id \ot Z_b X_a \bigr)$, whose terms are again
mutually orthogonal --- which is what makes the sum over the summed-out outcome free rather than a
factor $|\mB|$. One triangle inequality joins the two.
\end{proof}

\paragraph{Comments.} Formalized in \texttt{MIPRE/Foundations/Commutation.lean} (the fact) and
\texttt{MIPRE/Foundations/Pasting.lean} (the marginal step). This is the paper's
\texttt{lem:cool-closeness-fact} in the partition form its round-11 repair gives it, and the
partition form is the point: the consumer applies it to the $q$ fibres of an outcome map at once,
and the single-subset form would cost a factor $q$ there.

\begin{lemma}[Pasting, the two-measurement case]
\label{lem:pasting-updated}
\uses{def:cross-distance,lem:commutation-analysis,lem:cool-closeness-fact}
\lean{MIPRE.fibSum, MIPRE.isPVM_fibSum, MIPRE.sum_fiber, MIPRE.abs_sum_sum_le_sqrt,
  MIPRE.sum_comm4, MIPRE.sum_prod_eq, MIPRE.sum_prod_uniform, MIPRE.sum_prod_uniform_one,
  MIPRE.sum_weighted_add, MIPRE.sum_weighted_div, MIPRE.qform_conjTranspose,
  MIPRE.sum_sq_le_one_of_sum_eq_one, MIPRE.sum_snorm_sq_orth_le_one,
  MIPRE.sum_snorm_sq_povm_le_one, MIPRE.sum_snorm_sq_prod_le_one,
  MIPRE.sum_bornProb_le_fibSum, MIPRE.sum_xSqNorm_fibSum_le, MIPRE.sum_snorm_sq_comm_eq,
  MIPRE.pasteJ, MIPRE.sandOp_posSemidef, MIPRE.sum_sandOp, MIPRE.sum_snorm_sq_sandOp_le_one,
  MIPRE.sandOpG_posSemidef, MIPRE.sum_sandOpG, MIPRE.pasteJ_posSemidef, MIPRE.sum_pasteJ,
  MIPRE.abs_sigma_sub_cloud_le, MIPRE.abs_sand_sub_ord_le, MIPRE.sum_bornProb_ord_ge,
  MIPRE.sum_snorm_sq_comm_coarse_le, MIPRE.collisionTerm, MIPRE.collisionTerm_nonneg,
  MIPRE.strife_sub_cloud_eq, MIPRE.sum_snorm_sq_comm_fine_le,
  MIPRE.one_sub_sum_bornProb_pasteJ_le', MIPRE.one_sub_sum_bornProb_pasteJ_le,
  MIPRE.sum_xSqNorm_pasteJ_le, MIPRE.sum_collisionTerm_le}
\leanok
Let $\mu$ be a distribution on questions, $\ket{\psi}$ a unit vector, and suppose given, for each
question:
\begin{itemize}
\item a \emph{projective} measurement $\{A_{b,a}\}$ on $\C^{d_\alice}$ with outcomes pairs in
  $\mA \times \mB$;
\item a \emph{projective} measurement $\{R_b\}_{b \in \mA}$ on $\C^{d_\bob}$ --- in the
  application already the coarse-graining of the first family along its own outcome map, which is
  why only its projectivity appears;
\item a \emph{projective} measurement $\{G_g\}_{g \in \mC}$ on $\C^{d_\bob}$, a copy
  $\{G'_g\}$ of it on $\C^{d_\alice}$, and an outcome map $e : \mC \to \mB$.
\end{itemize}
Write $J_{b, a} = \sum_{g \,:\, e(g) = a} G_g R_b G_g$ for the pasted family, which is a POVM.
Suppose that
\begin{equation}
  \label{eq:pasting-hyp-bp}
  \bigl(\sum_a A_{b,a}\bigr) \ot \Id \simeq_\delta \Id \ot R_b\;, \qquad
  \bigl(\sum_b A_{b,a}\bigr) \ot \Id \simeq_\delta
    \Id \ot \sum_{g \,:\, e(g) = a} G_g\;, \qquad
  G'_g \ot \Id \simeq_\eta \Id \ot G_g\;,
\end{equation}
and that the \emph{collision term}
$\E \sum_b \sum_{g \neq g' \,:\, e(g) = e(g')}
  \bra{\psi} G'_{g'} \ot G_g R_b G_g \ket{\psi}$ is at most $\eps$. Then
\begin{equation}
  \label{eq:pasting-conclusion-bp}
  1 - \E \sum_{b, a} \bra{\psi} A_{b,a} \ot J_{b,a} \ket{\psi}
    \;\leq\; \frac{\delta}{2} + \sqrt{\frac{\delta}{2}}
      + \sqrt{32 \delta + 4\sqrt{\eta} + 2 \eps}\;,
\end{equation}
and hence $A_{b,a} \ot \Id \simeq_{\delta_{\mathrm{pasting}}} \Id \ot J_{b,a}$ with
$\delta_{\mathrm{pasting}}$ twice the right-hand side of~\eqref{eq:pasting-conclusion-bp}.

Finally, if the question is a pair $(z, y)$ with $y$ uniform and independent of $z$, every family
above depending on $z$ alone and only $e$ on $y$, and if any two distinct $g, g'$ satisfy
$e_z(y)(g) = e_z(y)(g')$ for at most an $\eps(z)$ fraction of the probes $y$, then the collision
term is at most the average of $\eps(z)$. The dependence on $z$ is what the consumer needs: for a
\emph{degenerate} line --- a direction that is zero, which the seeded test's diagonal branch does
sample --- the outcome map separates nothing at all, and what makes the average small is that such
lines are rare, not that the bound holds at every question.
\end{lemma}

\begin{proof}
\leanok
\uses{lem:commutation-analysis,lem:cool-closeness-fact,def:cross-distance}
The target is the agreement of $A$ with the sandwich. Two steps reach it. The ordered product
$G_g R_b$ carries all but $\delta/2 + \sqrt{\delta/2}$: its outer factor sums away against
Alice's second marginal by the second hypothesis, and what is left is Alice's first marginal
against $\Id - R_b$, whose Cauchy--Schwarz costs the square root of the first hypothesis. Passing
from the ordered product to the sandwich costs one more Cauchy--Schwarz against Bob's commutator
$[R_b, G_g]$, the front factor $A_{b, e(g)} \ot G_g$ being a mutually orthogonal family and so a
contraction.

So everything reduces to the \emph{fine-grained} commutator $\sum_{b, g} \norm{[R_b, G_g]
\ket{\psi}}^2$. Its coarse-grained counterpart, with $G$ replaced by its fibre sums, is at most
$16(\delta + \delta)$ by Lemma~\ref{lem:commutation-analysis} --- read in the mirror, with both
families on Bob's factor and Alice's joint measurement supplying the operator both products reach.
Expanding either commutator, the two diagonal terms are \emph{exactly one} because $R$ and $G$ are
projective, so each equals $2 - 2\Sigma$ for its own overlap $\Sigma$, and the two off-diagonal
terms are complex conjugates. Comparing the two overlaps is the only remaining work, and it passes
through the paper's cloud and strife: replacing Bob's outer factor by Alice's copy of it costs
$\sqrt\eta$ on each side --- and nothing depending on $|\mC|$, because one factor $R_b$
stays in front of the deviation, so the sum over $b$ is absorbed --- while the cloud and the strife
differ by exactly the collision term.

For the last part, the Born probabilities do not depend on the probe, so the probe average acts on
the collision indicator alone; what multiplies it is a POVM's total mass, which is one.
\end{proof}

\paragraph{Comments.} Formalized in \texttt{MIPRE/Foundations/Pasting.lean}. This is NW19's
Fact 4.35, the paper's \texttt{lem:pasting-updated} (\texttt{games.tex}), at $k = 2$; the paper's
reduction of general $k$ to $k = 2$ is an induction that adds nothing, and the consumer
(Lemma~\ref{lem:qld-pairs-of-lines}) needs $k = 2$ only.

\textbf{Two hypotheses of the paper's statement are absent, and both are absences rather than
gaps.} First, the inner family is assumed \emph{projective}; the paper allows a POVM there, and
then the two diagonal sums of the commutator expansion are only approximately one, which is what
its appeal to Fact 4.31 of~\cite{NW19} repairs. With projectivity they are exactly one and that
fact is not needed at all --- the consumer's inner family is a line measurement, which is
projective. Second, the cross-party self-consistency of $G$ at the \emph{fine} level is a
hypothesis here, where the paper derives it from the backwards consistency and Alice's
self-consistency and pays a further $\eps$ for the passage from the fibres to the fine family. The
consumer has the fine-level relation directly, as item 1 of
Lemma~\ref{lem:qld-expanded-lines}, so the derivation would be work spent to weaken a hypothesis
that is already met. Neither change touches the conclusion.

\textbf{The collision term is a hypothesis of the abstract lemma and a lemma of its own for the
product case.} That split is what keeps the analytic core free of the question distribution: the
core needs only that the off-diagonal cloud terms are small in total, and the product-distribution
bridge (\texttt{sum\_collisionTerm\_le}) is where uniformity of the probe and the separating
property of the outcome map are used.

\textbf{Constants.} $32\delta$ rather than the $16\delta$ of
Lemma~\ref{lem:commutation-analysis} because the two marginal hypotheses enter that lemma as a
single $\delta$, here instantiated at their sum. Nothing in the chain is a per-outcome estimate,
so no constant depends on $|\mC|$ or on $|\mB|$.

\subsection{Generalized Pauli observables}

\begin{definition}[Generalized Pauli observables in characteristic two]
\label{def:generalized-pauli}
\lean{MIPRE.sgn, MIPRE.sgn_add, MIPRE.star_sgn, MIPRE.norm_sgn, MIPRE.sum_sgn,
  MIPRE.Weyl.two_eq_zero, MIPRE.Weyl.add_self, MIPRE.Weyl.add_eq_zero_iff,
  MIPRE.Weyl.trDot, MIPRE.Weyl.trDot_comm, MIPRE.Weyl.trDot_add_left,
  MIPRE.Weyl.trDot_add_right, MIPRE.Weyl.trDotL, MIPRE.Weyl.wX, MIPRE.Weyl.wZ,
  MIPRE.Weyl.wX_mul_wX, MIPRE.Weyl.wZ_mul_wZ, MIPRE.Weyl.wX_mul_self,
  MIPRE.Weyl.wZ_mul_self, MIPRE.Weyl.wX_conjTranspose, MIPRE.Weyl.wZ_conjTranspose,
  MIPRE.Weyl.wX_isUnitary, MIPRE.Weyl.wZ_isUnitary, MIPRE.Weyl.wX_comm,
  MIPRE.Weyl.wZ_comm, MIPRE.Weyl.wX_mul_wZ, MIPRE.Weyl.wX_mul_wZ_of_ne,
  MIPRE.Weyl.wX_mul_wZ_of_eq, MIPRE.Weyl.trDotL_ne_zero, MIPRE.Weyl.sum_sgn_trDot,
  MIPRE.Weyl.sum_sgn_trDot_zero, MIPRE.Weyl.zProj, MIPRE.Weyl.sum_zProj,
  MIPRE.Weyl.wZ_eq_sum_zProj, MIPRE.Weyl.sum_sgn_smul_wZ,
  MIPRE.LowDegree.card_filter_ker_mul', MIPRE.Weyl.IsWeylFamily,
  MIPRE.Weyl.isWeylFamily_wZ, MIPRE.Weyl.isWeylFamily_wX, MIPRE.Weyl.proj,
  MIPRE.Weyl.proj_wZ, MIPRE.Weyl.proj_mul, MIPRE.Weyl.proj_mul_proj, MIPRE.Weyl.sum_proj,
  MIPRE.Weyl.eq_sum_proj, MIPRE.Weyl.proj_conjTranspose, MIPRE.Weyl.xProj,
  MIPRE.Weyl.sum_xProj, MIPRE.Weyl.wX_eq_sum_xProj, MIPRE.Weyl.xProj_mul_xProj,
  MIPRE.Weyl.xProj_conjTranspose, MIPRE.Weyl.dotF, MIPRE.Weyl.trDot_smul,
  MIPRE.Weyl.trMulL, MIPRE.Weyl.sum_sgn_linear, MIPRE.Weyl.sum_sgn_trMul,
  MIPRE.Weyl.syn, MIPRE.Weyl.sum_sgn_smul_line}
\leanok
Let $q = 2^k$ be an admissible field size and let $n$ be a
finite index set, so that $\F_q^n$ indexes a basis of $(\C^q)^{\ot n}$. Write
$\fieldtr : \F_q \to \F_2$ for the field trace and, for $a, b \in \F_q^n$,
\begin{equation*}
  \langle a, b \rangle \;=\; \fieldtr\Bigl(\sum_{l \in n} a_l b_l\Bigr) \;\in\; \F_2\;.
\end{equation*}
The \emph{generalized Pauli}, or \emph{Weyl}, \emph{observables} are
\begin{equation*}
  \qp^X(a) \;=\; \sum_{i \in \F_q^n} \ket{i + a}\bra{i}\;, \qquad
  \qp^Z(b) \;=\; \sum_{i \in \F_q^n} (-1)^{\langle b, i \rangle}\, \ket{i}\bra{i}\;.
\end{equation*}
They satisfy, for all $a, a', b, b' \in \F_q^n$:
\begin{enumerate}
\item $\qp^W(a) \qp^W(a') = \qp^W(a + a')$, so $\qp^W(0) = \Id$ and each family commutes;
\item $\qp^W(a)^\dagger = \qp^W(a)$ and $\qp^W(a)^2 = \Id$: each is a self-adjoint involution,
  hence a $\pm 1$-observable and a unitary;
\item (\emph{twisted commutation})
  $\qp^X(a)\, \qp^Z(b) = (-1)^{\langle a, b \rangle}\, \qp^Z(b)\, \qp^X(a)$; so the two commute
  when $\langle a, b\rangle = 0$ and anticommute when it is $1$.
\end{enumerate}
For $W \in \{X, Z\}$ define the \emph{spectral projectors} of the family by Fourier averaging,
\begin{equation}
  \label{eq:pauli-inversion-bp}
  \qp^W_e \;=\; \E_{a \in \F_q^n} (-1)^{\langle a, e \rangle}\, \qp^W(a)\;.
\end{equation}
For $W = Z$ these are the projectors onto the computational basis vectors. They form a projective
measurement and invert~\eqref{eq:pauli-inversion-bp}:
\begin{equation}
  \label{eq:pauli-obs-proj-bp}
  \qp^W_e \qp^W_{e'} = \delta_{e, e'}\, \qp^W_e\;, \qquad
  \sum_e \qp^W_e = \Id\;, \qquad
  \qp^W(a) \;=\; \sum_{e} (-1)^{\langle a, e\rangle}\, \qp^W_e\;.
\end{equation}
Finally, for $v \in \F_q^n$ and $a \in \F_q$ write $v \cdot e = \sum_l v_l e_l \in \F_q$ for the
pairing before the trace, and define the \emph{syndrome projector}
$\qp^W_{[\,\cdot\, v = a]} = \sum_{e \,:\, e \cdot v = a} \qp^W_e$. Averaging the family along the
line $r \mapsto r v$ recovers it:
\begin{equation}
  \label{eq:pauli-line-bp}
  \E_{r \in \F_q} (-1)^{\fieldtr(a r)}\, \qp^W(r v) \;=\; \qp^W_{[\,\cdot\, v = a]}\;.
\end{equation}
\end{definition}

\paragraph{Source.} The pinned paper~\cite{AuditSync26}, \texttt{preliminaries.tex},
\texttt{sec:generalized-pauli}: \texttt{eq:pauli-fp} and its $\F_q$ analogue,
\texttt{eq:twisted-fq}, \texttt{eq:pauli-product-power}, \texttt{eq:pauli-obs-proj} and
\texttt{eq:pauli-inversion-0}. The cancellation behind the inversion is the paper's
\texttt{lem:cancellation} (Fact 3.2 of~\cite{NW19}).

\paragraph{Comments.} The formalization is \texttt{MIPRE/Foundations/Weyl.lean}. It is the half of
the Pauli basis test's infrastructure that was missing: \texttt{MIPRE/LCS/Pauli.lean} is the
$2 \times 2$ Pauli matrices for the Magic Square, and what the expansion stage of
Remark~\ref{rem:qld-architecture} and its swap isometry both run on is this $\F_q$ Weyl system.
Three points, two of them restrictions stated so that no later statement is read as more general
than it is.

\textbf{Characteristic two only.} The paper gives the construction over $\F_p$ with phases
$\omega = e^{2\pi i/p}$ and then specializes; this is the $p = 2$ case throughout, where
$\omega = -1$ and every phase is the real sign \texttt{MIPRE.sgn}. Nothing is lost here: the test's
parameters are admissible by hypothesis, which means $q = 2^k$, and no consumer instantiates it
elsewhere. What it buys is that the operators are self-adjoint \emph{involutions} rather than
order-$p$ unitaries --- so they are genuine $\pm 1$-observables, and item 2 above is where the
paper needs $(\qp^W(a))^p = \Id$.

\textbf{One matrix, not an iterated Kronecker product.} $(\C^q)^{\ot n}$ is presented as matrices
indexed by the functions $n \to \F_q$, which \emph{is} $\F_q^n$, so $\qp^W(a)$ is a single matrix
and the $n$-register statements are the one-register statements with no associativity
bookkeeping. The index set $n$ is an arbitrary finite type rather than $\{1, \ldots, n\}$, which is
what lets the test use it at $n = \{0,1\}^m$ --- the indexing of $\F_q^M$, $M = 2^m$, that
Definition~\ref{def:ld-encoding} already uses.

\textbf{The cancellation lemma needed a generalization, not a port.} The identity behind the
inversion in~\eqref{eq:pauli-obs-proj-bp} is that a nontrivial sign character of $\F_q^n$ sums to
zero, and that is the index-two kernel of a nonzero $\F_2$-linear functional. The repository had
that count only for a functional on a field \emph{extension}
(\texttt{card\_filter\_ker\_mul}); the proof uses nothing but the module structure and a vector
off the kernel, so it is now stated for any finite $\F$-vector space
(\texttt{card\_filter\_ker\_mul'}), with the field-extension form as the instance. The paper's own
correction note on \texttt{lem:cancellation} warns that the ambient space must be a subspace over the
\emph{base} field and that the statement is false over a subfield; here the space is all of
$\F_q^n$ over $\F_2$, so there is nothing to check.

\textbf{The spectral theory is proved once, for both families.} The only properties of a family
that~\eqref{eq:pauli-obs-proj-bp} and~\eqref{eq:pauli-line-bp} use are the three of
\texttt{IsWeylFamily}: it represents the additive group, its operators are self-adjoint, and the
identity sits at $0$. So the projectors are built by~\eqref{eq:pauli-inversion-bp} from any such
family and the statements hold for both, which means the paper's Fourier transform
(\texttt{eq:fourier-f}) conjugating $\qp^Z$ into $\qp^X$ is \emph{not needed} to reach the
$X$-side projectors. That is worth avoiding: the transform carries a $q^{-1/2}$, and a square root
in matrix entries costs far more in Lean than it does on paper.

Two of the identities turn out to need even less. The expansion of an operator in its own
projectors, and the line average~\eqref{eq:pauli-line-bp}, use \emph{no} group law at all --- they
are Fourier inversion on the index group and hold for an arbitrary family. What the group law buys
is that the projectors are mutually \emph{orthogonal}, which is what makes them a measurement
rather than a resolution.

Equation~\eqref{eq:pauli-line-bp} is the identity Lemma~\ref{lem:qld-expanded-points} averages
to, at $v = \ind_m(u)$.

\begin{definition}[The maximally entangled state of the Weyl system]
\label{def:weyl-epr}
\uses{def:generalized-pauli,def:state-distance}
\lean{MIPRE.Weyl.eprScale, MIPRE.Weyl.epr, MIPRE.Weyl.epr_same, MIPRE.Weyl.epr_ne,
  MIPRE.Weyl.eprScale_sq, MIPRE.Weyl.epr_unit, MIPRE.Weyl.stateVec_entry,
  MIPRE.Weyl.stateVec_epr, MIPRE.Weyl.wX_transpose, MIPRE.Weyl.wZ_transpose,
  MIPRE.Weyl.stateVec_epr_wX, MIPRE.Weyl.stateVec_epr_wZ, MIPRE.Weyl.stateVec_epr_proj,
  MIPRE.sgn_sum, MIPRE.Weyl.isPVM_syn, MIPRE.Weyl.stateVec_epr_syn,
  MIPRE.Weyl.synOf, MIPRE.Weyl.syn_eq_synOf, MIPRE.Weyl.sum_synOf,
  MIPRE.Weyl.isPVM_synOf, MIPRE.Weyl.stateVec_epr_synOf}
\leanok
On two copies of $(\C^q)^{\ot n}$,
\begin{equation*}
  \ket{\EPR_q}^{\ot n} \;=\; \frac{1}{\sqrt{q^n}} \sum_{a \in \F_q^n} \ket{a} \ot \ket{a}\;,
\end{equation*}
a unit vector. For any operator $M$ on one copy,
\begin{equation}
  \label{eq:epr-transport-bp}
  (M \ot \Id) \ket{\EPR_q}^{\ot n} \;=\; (\Id \ot M^{\mathsf T}) \ket{\EPR_q}^{\ot n}\;,
\end{equation}
and the Weyl operators and their spectral projectors are \emph{symmetric}, so for them the
transpose is invisible: the two parties' operators act identically on the state.

Both facts about the syndrome projectors $\qp^W_{[\,\cdot\, v = a]}$ --- that they are a projective
measurement, and that they transport across the state --- hold for the coarse-graining of the
eigenbasis measurement along a level set of an \emph{arbitrary} label, the syndrome being the case
of the label $e \mapsto \langle e, v\rangle$; and coarse-graining such a family again gives the
family of the composite label.
\end{definition}

\paragraph{Comments.} Formalized in \texttt{MIPRE/Foundations/WeylEPR.lean}. The stabilizer
relation of~\eqref{eq:epr-transport-bp} at $M = \qp^W(a)$ is what the expansion stage of
Remark~\ref{rem:qld-architecture} runs on: its expanded state adjoins $\ket{\EPR_q}^{\ot M}$ to
each player, and the sign $(-1)^{\gamma(\omega)}$ that the strategy's own observables carry is
cancelled against one the \emph{ancilla} observables carry, for which the two ancilla halves must
be tied together.

This is the second dividend of characteristic two, after the operators being involutions. Over
$\F_p$ for odd $p$ the transpose of $\qp^Z(b)$ is $\qp^Z(b)$ but its conjugate is $\qp^Z(-b)$, and
the bookkeeping reappears; here every entry is real and the transpose is the identity map on the
whole family, projectors included (\texttt{stateVec\_epr\_proj}, which needs only that the family
is symmetric).

\effortEasy

\begin{lemma}[Qudits are qubits, in characteristic two]
\label{lem:pauli-binary}
\uses{def:generalized-pauli,def:weyl-epr,lem:self-dual-basis}
\lean{MIPRE.Weyl.binEquiv, MIPRE.Weyl.binEquiv_apply, MIPRE.Weyl.binEquiv_add,
  MIPRE.Weyl.binEquiv_injective, MIPRE.Weyl.binEquiv_eq_iff, MIPRE.Weyl.trDot_binEquiv,
  MIPRE.Weyl.wX_binEquiv, MIPRE.Weyl.wZ_binEquiv, MIPRE.Weyl.card_binEquiv,
  MIPRE.Weyl.proj_binEquiv, MIPRE.Weyl.proj_wZ_binEquiv, MIPRE.Weyl.proj_wX_binEquiv,
  MIPRE.Weyl.eprScale_binEquiv, MIPRE.Weyl.epr_binEquiv, MIPRE.Weyl.pauli_binary,
  MIPRE.Weyl.exists_binEquiv_pauli_binary, MIPRE.Weyl.trDot_two, MIPRE.Weyl.sgn_trDot_two,
  MIPRE.Weyl.wZ_two_diag, MIPRE.Weyl.wX_two_apply}
\leanok
Let $q = 2^k$ be an admissible field size, let $L$ be a finite index set, and let
$\{e_1, \ldots, e_k\}$ be a self-dual basis of $\F_q$ over $\F_2$, with coordinate map
$\kappa$. Then the relabelling
\begin{equation*}
  \phi : \F_q^L \to \F_2^{L \times k}\;, \qquad \phi(u)_{i j} \;=\; \kappa(u_i)_j\;,
\end{equation*}
is a bijection, and conjugating by it carries the whole Weyl system over $\F_q$ on $L$ registers
to the Weyl system over $\F_2$ on $L k$ registers: the form $\langle \cdot, \cdot \rangle$, the
two operator families $\qp^X$ and $\qp^Z$, their spectral projectors, and
$\ket{\EPR_q}^{\ot L} \mapsto \ket{\EPR_2}^{\ot L k}$. Such a basis exists whenever $k$ is odd.
\end{lemma}

\begin{proof}
\leanok
\uses{def:generalized-pauli,def:weyl-epr,lem:self-dual-basis}
Four steps, one per piece, and only the first uses self-duality. The form transports because
$\fieldtr(x y) = \kappa(x) \cdot \kappa(y)$, which is item 2 of the paper's
\texttt{lem:downsize\_field} and is exactly self-duality
(\texttt{MIPRE.LowDegree.trace\_mul\_eq\_dot}); summing over the $L$ registers is
\texttt{trDot\_binEquiv}. $\qp^Z$ transports because its entries are the characters of the form,
so it transports as soon as the form does. $\qp^X$ transports because $\phi$ is
\emph{additive}, which is all a shift knows about --- no self-duality. The projectors transport
because the Fourier average of~\eqref{eq:pauli-inversion-bp} reindexes along $\phi$ and the
normalization is one cardinality counted on either side. The state transports because both sides
are the normalized diagonal and a relabelling preserves the diagonal. Existence of the basis is
the existence half of Lemma~\ref{lem:self-dual-basis}.
\end{proof}

\paragraph{Comments.} Formalized in \texttt{MIPRE/Foundations/WeylBinary.lean}; this is the
paper's \texttt{lem:pauli-binary} (\texttt{preliminaries.tex}), which is what makes the Pauli
basis test a self-test for \emph{qubit} Pauli observables and $\ket{\EPR_2}$, as its
\texttt{cor:pauli-binary} states.

\textbf{A relabelling, not an isometry.} The paper builds an isometry
$\theta : \C^q \to (\C^2)^{\ot k}$ and states the conclusion as
$\qp^W_u = \phi^\dagger (\bigotimes_{i,j} \sigma^W_{u_{ij}}) \phi$. Here both systems are
matrices indexed by a finite set --- $\F_q^L$ on one side, $\F_2^{L \times k}$ on the other ---
and $\phi$ is a bijection \emph{of those index sets}, so conjugating by the corresponding
permutation unitary is \texttt{Matrix.submatrix} along $\phi$ and no unitary has to be
constructed. That is the same simplification as presenting $(\C^q)^{\ot n}$ as one matrix rather
than an iterated Kronecker product, and it is why the four items are four short calculations.

\textbf{What ``these are qubit Paulis'' means here.} The target is the Weyl system at
$\F = \F_2$ and index set $L \times k$, whose matrices are indexed by
$\F_2^{L \times k}$ --- the computational basis of $(\C^2)^{\ot L k}$. That the operators really
are the tensor products of single-qubit Paulis is then a statement \emph{inside} that system, and
it is checked rather than read off the encoding: the form is the plain dot product
(\texttt{trDot\_two}, the field trace of a one-dimensional extension being the identity), its
character is a product over the qubits (\texttt{sgn\_trDot\_two}), and each operator's entries
factor accordingly (\texttt{wZ\_two\_diag} on the diagonal, \texttt{wX\_two\_apply} everywhere).
The paper's $\bigotimes_{i=1}^L \bigotimes_{j=1}^k \sigma^W_{u_{ij}}$ is that factorization.

\textbf{The oddness of $k$ is inherited, not needed.} Every statement above holds for any
self-dual basis, given as a hypothesis; only the existence clause asks $k$ to be odd, and that is
the restriction of Lemma~\ref{lem:self-dual-basis}'s formalized half, not of this lemma. The
paper's admissible sizes satisfy it.

\begin{definition}[Inconsistency]
\label{def:inconsistency}
\uses{def:game}
\lean{MIPRE.inconsistency}
\leanok
Let $\mu$ be a distribution on a finite set $\mX$, let
$\ket{\psi} \in \C^{d_\alice} \ot \C^{d_\bob}$, and let $\{M^x_a\}$, $\{N^x_a\}$ be
families of POVMs on the two factors with a common question set $\mX$ and a common
answer set. Their \emph{inconsistency} on $\ket{\psi}$ is
\begin{equation*}
  \E_{x \sim \mu} \sum_{a \neq b} \bra{\psi} M^x_a \ot N^x_b \ket{\psi},
\end{equation*}
the probability that the two measurements return different answers. We say $M$ and $N$
are \emph{$\delta$-consistent} if their inconsistency is at most $\delta$.
\end{definition}

\paragraph{Comments.} This is Definition 4.8 of~\cite{JNVWY21qld}, in the two-space
form used by Theorem~\ref{thm:lidt-soundness}. It is a different notion from the POVM
distance of Definition~\ref{def:distance}: consistency compares two families across the
two factors of the state, whereas $\approx_\delta$ compares them as operator families.

\begin{remark}[The POVM distance and the state]
\label{rem:distance-state}
\uses{def:distance,def:sync-strategy,def:sync-value}
Definition~\ref{def:distance} is \emph{state-independent}, whereas the companion paper's
$\approx_\delta$ (\texttt{games.tex}) is the state-dependent quantity $\E_{x \sim \mu}
\sum_a \norm{(M^x_a - N^x_a)\ket{\psi}}^2$. For a \emph{synchronous} strategy the two agree.
The state is then the maximally entangled state on $\C^d \ot \C^d$, and for it
$\bra{\psi} (A^\dagger A) \ot I \ket{\psi} = \tr(A^\dagger A)/d = \norm{A}_{hs}^2$, so
Definition~\ref{def:distance} \emph{is} the state-dependent distance at the tracial state.
Every formalized consumer of Definition~\ref{def:distance} is synchronous, and the soundness
of oracularization (Lemma~\ref{lem:oracular-soundness}) reads it that way. For a general
bipartite strategy the two notions differ, so a chapter~3 consumer that is not synchronous
will need the state-dependent notion added rather than this one reused.
\end{remark}

\begin{lemma}[Closeness from two disagreements]
\label{lem:hs-closeness}
\uses{def:hs-norm,def:distance}
\lean{MIPRE.sum_hsNormSq_sub_le}
\leanok
Let $\{O_i\}_{i \in I}$, $\{C_i\}_{i \in I}$, $\{D_i\}_{i \in I}$ be finite families of
projections on $\C^{d \times d}$ with $\sum_i \tau(O_i) = 1$ and $\sum_i \tau(C_i D_i) = 1$.
Then
\begin{equation*}
  \sum_i \norm{C_i D_i - O_i}_{hs}^2 \;\leq\; 3 \Bigl( \sum_i \tau\bigl(O_i (1 - D_i)\bigr)
    + \sum_i \tau\bigl(O_i (1 - C_i)\bigr) \Bigr).
\end{equation*}
\end{lemma}
\begin{proof}
\leanok
Write $X_i = (1 - C_i) O_i$ and $Y_i = (1 - D_i) O_i$, so that $C_i O_i = O_i - X_i$ and
$D_i O_i = O_i - Y_i$. Since $1 - C_i$ and $1 - D_i$ are projections,
$\norm{X_i}_{hs}^2 = \tau(O_i (1 - C_i))$ and $\norm{Y_i}_{hs}^2 = \tau(O_i (1 - D_i))$; write
$\beta$ and $\alpha$ for the two sums, which are therefore nonnegative. Expanding the
left-hand side and using $\sum_i \tau(C_i D_i) = \sum_i \tau(O_i) = 1$,
\begin{equation*}
  \sum_i \norm{C_i D_i - O_i}_{hs}^2 \;=\; 2 - 2 \sum_i \tau(O_i D_i C_i).
\end{equation*}
Since $O_i$ is a projection, $\tau(O_i D_i C_i) = \langle D_i O_i, C_i O_i \rangle_{hs} =
\langle O_i - Y_i, O_i - X_i \rangle_{hs}$, which expands to $\tau(O_i) -
\tau(O_i(1 - C_i)) - \tau(O_i(1 - D_i)) + \langle Y_i, X_i \rangle_{hs}$. Summing over $i$,
$\sum_i \tau(O_i D_i C_i) = 1 - \alpha - \beta + \sum_i \langle Y_i, X_i \rangle_{hs}$, and
Cauchy--Schwarz followed by the arithmetic--geometric mean inequality gives
$\bigl| \sum_i \langle Y_i, X_i \rangle_{hs} \bigr| \leq \sqrt{\alpha}\sqrt{\beta} \leq
(\alpha + \beta)/2$. Hence $\sum_i \norm{C_i D_i - O_i}_{hs}^2 \leq 2(\alpha + \beta) +
(\alpha + \beta)$.
\end{proof}

\begin{lemma}[Close measurements have close values]
\label{lem:close-measurements-close-values}
\uses{def:hs-norm,def:distance}
\lean{MIPRE.sum_ntr_mul_ge}
\leanok
Let $\{O_i\}_{i \in I}$, $\{C_i\}_{i \in I}$, $\{D_i\}_{i \in I}$ be finite families of
projections on $\C^{d \times d}$ with $\sum_i \tau(O_i) = 1$, and let
$\mathrm{acc} \subseteq I$. Then
\begin{equation*}
  \sum_{i \in \mathrm{acc}} \tau(O_i)
    \;-\; 2 \sqrt{\sum_i \norm{C_i D_i - O_i}_{hs}^2}
    \;\leq\; \sum_{i \in \mathrm{acc}} \tau(C_i D_i).
\end{equation*}
\end{lemma}
\begin{proof}
\leanok
This is the measurement-level form of \cite[Fact 4.31]{NW19}, in the one direction a soundness
argument needs and in the synchronous setting, where $\tau((C_iD_i)^\dagger C_iD_i) =
\tau(C_iD_i)$ makes the products a measurement on the nose. Put $w_i = C_iD_i - O_i$. Then
$\tau(C_iD_i) = \norm{C_iD_i}_{hs}^2 = \norm{O_i + w_i}_{hs}^2 = \tau(O_i) +
2\langle O_i, w_i\rangle_{hs} + \norm{w_i}_{hs}^2$. Summing over $\mathrm{acc}$ and dropping
the last term, which is nonnegative,
\begin{equation*}
  \sum_{i \in \mathrm{acc}} \tau(C_iD_i) \;\geq\; \sum_{i \in \mathrm{acc}} \tau(O_i)
    - 2 \Bigl| \sum_{i \in \mathrm{acc}} \langle O_i, w_i \rangle_{hs} \Bigr|,
\end{equation*}
and Cauchy--Schwarz over the index bounds the last modulus by
$\sqrt{\sum_{i \in \mathrm{acc}} \tau(O_i)} \cdot \sqrt{\sum_{i \in \mathrm{acc}}
\norm{w_i}_{hs}^2} \leq 1 \cdot \sqrt{\sum_i \norm{w_i}_{hs}^2}$.
\end{proof}

\paragraph{Comments.} The square root is spent here and nowhere else in the soundness of
oracularization; everything before it, including Lemma~\ref{lem:hs-closeness}, is linear.
That is what repair~1 of Remark~\ref{rem:oracularization-repairs} is about.

\subsection{Computability conventions}
\label{sec:computability-conventions}

Turing machines are represented by their G\"odel numbers; in Lean the intended model is
Mathlib's \texttt{Nat.Partrec.Code}, with $\phi_e$ the partial function computed by
code $e$ and ``$e$ halts on empty input'' meaning $(\texttt{e.eval } 0).\texttt{Dom}$.
Running times are counted as in~\cite{MNY}; polynomial-time computability of a function
of strings and integers means time polynomial in the lengths of the string inputs and
in $\log$ of the integer inputs, unless an integer is explicitly given in unary.

\begin{definition}[Recursively enumerable]
\label{def:re}
\lean{MIPRE.IsRE, MIPRE.isRE_iff}
\leanok
A language $L \subseteq \{0,1\}^*$ is in $\RE$ if it is the domain (equivalently, the
range) of a partial computable function. The \emph{halting problem} --- the set of $e$
such that $\phi_e$ halts on the empty input --- is $\RE$-complete under many-one
reductions.
\end{definition}

In Lean (\texttt{MIPRE.IsRE}) membership in $L$ is Mathlib's \texttt{REPred}, the domain
of a partial computable function; equivalently (\texttt{MIPRE.isRE\_iff}), $L$ is the
halting set of a well-scoped program of the ambient model
(Section~\ref{sec:rr-computability}) on encoded strings, which is the form in which the
compressibility criterion (Lemma~\ref{lem:compressible-criterion}) consumes recursive
enumerability. The completeness of the halting problem is
Theorem~\ref{thm:halting-undecidable}, not part of the definition.

\begin{definition}[Succinct description]
\label{def:succinct}
\lean{MIPRE.Cost.IsSuccinctDesc, MIPRE.Cost.bitQueryAnswer}
\leanok
Following~\cite{MNY}, a pair $(e, n)$ of a Turing machine $e$ and integer $n$ is a
\emph{succinct description} of the string $x \in \{0,1\}^*$ if $\abs{x} \leq 2^n$, the
runtime of $e$ on input $m$ is at most $(n+1)(\abs{m}+1)^2$, and $\phi_e(m)$ equals the
$m$-th bit of $x$ for $m < \abs{x}$, and $2$ for $m \geq \abs{x}$. (\cite{MNY} require
runtime at most $n$ on every input; in a model that charges for reading its input this is
vacuous for $\abs{m} > n$, and the polynomial slack in $\abs{m}$ is what the ambient
model's bit-indexing program achieves. Any fixed polynomial in $n$ and $\abs{m}$ would
do, provided the same one is used in Theorem~\ref{thm:compression}.)
\end{definition}

\begin{definition}[Game description]
\label{def:game-description}
\uses{def:sync-game}
\lean{HaltingGameValue.GameData, HaltingGameValue.GameData.toGame,
  HaltingGameValue.GameData.game}
\leanok
A \emph{game description} is a tuple $g = (n_\mX, n_\mA, w, \mathrm{acc})$ of two
integers, a finite list $w$ of weighted question pairs, and a finite list
$\mathrm{acc}$ of accepted answer tuples. It determines a synchronous game
$\game_g$ on question alphabet $\{0, \ldots, n_\mX\}$ and answer alphabet $\{0, \ldots,
n_\mA\}$: the question distribution is $w$ normalized (a fixed point mass if $w$ has
total weight zero), and the predicate accepts exactly the tuples listed in
$\mathrm{acc}$, except that unequal answers to equal questions always lose. Game
descriptions are first-order data and carry a canonical G\"odel numbering
(\texttt{GameData} in the Lean file).
\end{definition}

\subsection{Normal form verifiers}
\label{sec:normal-form-verifiers}

Games produced in the proof are presented by pairs of Turing machines: a
\emph{sampler}, generating questions by a rigid linear-algebraic procedure that
introspection can certify, and a \emph{decider}, checking answers.

\begin{definition}[Sampler]
\label{def:sampler}
\uses{def:cl-function,lem:cl-structure}
\lean{MIPRE.CL.Sampler, MIPRE.CL.Sampler.Query, MIPRE.CL.Sampler.dist,
  MIPRE.CL.Sampler.TimeBound}
\leanok
An \emph{$\ell$-level sampler} (with field size $2$) is a Turing machine $\sampler$ together
with, for each $n$, a dimension $s(n)$ and two $\ell$-level CL functions $L^\alice, L^\bob$ on
$V = \F_2^{s(n)}$ given with exact decomposition data (Lemma~\ref{lem:cl-structure}) such
that, writing vectors and register subspaces of $V$ as bit strings of length $s(n)$
(coordinates, respectively indicator vectors), for all $w \in \{\alice, \bob\}$, $j \in \{1,
\ldots, \ell\}$ and bit strings $z, u, y$ of length $s(n)$:
\begin{itemize}
\item on input $(n, \mathtt{dimension})$, $\sampler$ outputs $s(n)$;
\item on input $(n, w, \mathtt{marginal}, j, z)$, it outputs the marginal $L^w_{\leq j}(z)$;
\item on input $(n, w, \mathtt{linear}, j, u, y)$, it outputs $L^w_{j,u}(y)$, provided
  $u \in L^w_{<j}(V)$;
\item on input $(n, w, \mathtt{factor}, j, u)$, it outputs the factor space $V^w_{j,u}$,
  provided $u \in L^w_{<j}(V)$;
\end{itemize}
and $\sampler$ halts on every input of the form $(n, \ldots)$, well formed or not; nothing is
required of its outputs on other inputs. The \emph{distribution of $\sampler$ at index $n$}
is the CL distribution $\mu_n$ of $(L^\alice(x), L^\bob(x))$ for uniform $x \in V$
(Definition~\ref{def:cl-function}). In the paper's Turing-machine model,
$\TIME_\sampler(n)$ denotes the absolute supremum of the running times over all inputs
of the form $(n, \ldots)$; totality alone does not make that supremum finite.
\end{definition}

\paragraph{Comments.} This follows the query interface of the paper's
\texttt{def:sampler} (\texttt{linear.tex})~\cite{AuditSync26} at
field size $2$, the only size the normal form verifiers of this blueprint use; the paper's
samplers over $\F_{2^t}$ reach this form by downsizing (Section~\ref{sec:ps-cl}). Two
interface details matter here. The correctness clauses for $\mathtt{linear}$ and $\mathtt{factor}$
are imposed only for prefixes in the range of $L^w_{<j}$, as the paper's amended definition
states, and are well defined there by Lemma~\ref{lem:cl-structure}. The paper additionally
requires bounded-prefix format and length checks: after computing $s(n)$, a sampler
inspects only the prescribed prefix of each auxiliary input, rejecting with a fixed default
when the format fails. Together with totality, this bounds the paper's absolute
$\TIME_\sampler(n)$ uniformly over even malformed or overlong inputs. A total program
that scans an ignored, arbitrarily long suffix would not satisfy that source contract.

In Lean the machine is a program of the ambient cost model
(Section~\ref{sec:computability-conventions}), its well-formed input is the pair of $n$
and a query (\texttt{Sampler.Query}), and \texttt{Sampler.TimeBoundAt n T k} is a
\emph{relative cost bound}: it requires halting within
$T\cdot(\abs{d}+1)^k$ on every input $(n,d)$, well formed or not. It is a separate
predicate, not a bound $\TIME_\sampler(n)\leq T$ on the absolute supremum. In particular,
neither this predicate nor the totality field supplies the paper's bounded-prefix contract.

Two departures from the paper are forced here, and both were found by trying to prove the
bound for a concrete machine rather than by inspection. First, the paper's bound is
\emph{uniform} over all inputs, and that reading is not available in the ambient model: a
value is inspected in place one node at a time, but a loop carries the rest of a list to its
next iteration only by copying it, so no program walks an unbounded list within a cost
independent of its length, and a uniform bound is met by no sampler or decider that checks
the length of a string against an unbounded $s(n)$. The running time is therefore bounded by a
polynomial in the size of the input. Second, the \emph{degree} $k$ of that polynomial has to
be carried alongside the coefficient $T$, because every decider of the pipeline runs another
program through the universal machine, whose overhead is a polynomial $Q$ of degree above one
(\texttt{Cost.Universal.univPoly}): a simulated run costing $T (\abs{d}+1)^k$ becomes one
costing about $Q(T (\abs{d}+1)^k)$, of degree $k \cdot \deg Q$. No fixed degree is closed
under the composition the deciders perform; a bounded degree is. Collapsing the two into one
parameter --- reading the bound as $T (\abs{d}+1)^T$ --- would also be closed, but then
$\lambda$-boundedness would allow degree $n^\lambda$ in the input size, which is not a
polynomial-time condition at all and would make Theorem~\ref{thm:compression} assume
compression of verifiers its proof cannot compress.

\begin{definition}[Decider]
\label{def:decider}
\lean{MIPRE.Decider, MIPRE.Decider.Accepts, MIPRE.Decider.TimeBound}
\leanok
A \emph{decider} is a $5$-input Turing machine $\decider$, whose input is interpreted
as $(n, x, y, a, b)$ with $n$ an integer index and $x, y, a, b$ strings. It \emph{accepts}
if it halts with output $1$, and \emph{rejects} otherwise (output $0$, or a timeout).
In the paper, $\TIME_\decider(n)$ denotes the absolute supremum of its running times
over all inputs with index $n$, possibly infinite. In Lean,
\texttt{Decider.TimeBoundAt n T k} instead requires halting within cost
$T \cdot (\abs{d} + 1)^k$ on every input $(n, d)$. This relative bound implies
totality at index $n$, but need not imply finite absolute $\TIME_\decider(n)$, for the
reasons given after Definition~\ref{def:sampler}.
\end{definition}

\begin{definition}[Normal form verifier in the ambient model]
\label{def:normal-verifier}
\uses{def:sampler,def:decider,def:game,def:sync-game}
\lean{MIPRE.Verifier, MIPRE.Verifier.game, MIPRE.Verifier.Answers, MIPRE.Verifier.size,
  MIPRE.Verifier.IsSynchronousAt, MIPRE.Verifier.syncGame}
\leanok
An ambient \emph{normal form verifier} is a pair $\verifier=(\sampler,\decider)$ of
a sampler over $\F_2$ and a decider satisfying the extensional condition
\begin{equation*}
  \decider(n,x,y,a,b)\text{ accepts}\;\Rightarrow\;
  \abs{x}=\abs{y}=s(n).
\end{equation*}
This is \texttt{Verifier.accepts\_length}; no particular sequence of preliminary checks
is part of the structure. For each index $n$ and independently chosen answer cut
$T\in\N$, the game $\verifier_{n,T}$ (\texttt{Verifier.game n T}) has question sets
$\mX=\mY=\{0,1\}^{s(n)}$, answer sets consisting of the bit strings of length at most
$T$, the sampler distribution at index $n$, and decision predicate equal to the decider's
acceptance predicate. Its definition imposes no time-bound hypothesis. We write
$\verifier_n$ when the answer cut is specified by the surrounding statement.

The condition \texttt{Verifier.IsSynchronousAt n} says that the decider rejects unequal
answers to equal questions at index $n$. Under this condition,
\texttt{Verifier.syncGame n T} is the same game as a synchronous game, so its synchronous
value is defined (Remark~\ref{rem:sync-invariant}). The description length is
$\abs{\verifier} = \max\{\abs{\sampler}, \abs{\decider}\}$, and the number of levels of
$\verifier$ is that of $\sampler$.
\end{definition}

\begin{remark}[Comparison with the source normal-verifier contract]
\label{rem:source-normal-verifier}
\ledgernode{1.1.3}
\uses{def:normal-verifier,rem:runtime-contract-bridge}
The paper's \texttt{def:normal-ver} requires an operational first step: compute $s(n)$
and check both question lengths before continuing. Its \texttt{def:normal-game} uses
answer strings of length at most the finite absolute time $\TIME_\decider(n)$
(\texttt{games.tex}~\cite{AuditSync26}). The dimension check yields
$s(n)\leq\TIME_\decider(n)$ in that machine model. These source requirements are not
the checked ambient definition above.

Relating executions in the two models and identifying the source answer cut are part of
the pending normalization in Remark~\ref{rem:runtime-contract-bridge}. An ambient
relative-cost coefficient need not be the source absolute time, and increasing an answer
cut may change the value. In the ambient model acceptance is halting with output $1$;
the relative time bound ensures totality with input-dependent budgets. The dimension bound
needed by the pipeline is required independently in Definition~\ref{def:lambda-bounded}.
\end{remark}

\begin{remark}[The synchronization invariant]
\label{rem:sync-invariant}
\uses{def:normal-verifier,def:sync-game,thm:almost-sync}
An earlier arrangement of the pipeline of Section~\ref{ch:proof-structure} asked two
properties of the games it produces, neither of which follows from symmetrizing the
players' roles. The pipeline no longer needs either, its soundness being carried in
$\valstar$ throughout (Remark~\ref{rem:bipartite-route}); the two are recorded here
because they are what a synchronous route would have to establish, and because the second
is the reason there is no synchronous route:
\begin{enumerate}
\item \emph{Synchronicity}: equal questions reject unequal answers, so that $\synval$ is
  defined at all. For a game given by a game description this holds by fiat
  (Definition~\ref{def:game-description}); for a game given by a normal form verifier it
  is a condition on the decider. It would be needed wherever $\synval$ occurs; direct
  repetition preserves it, since two unequal answer tuples to equal question tuples differ
  in some coordinate, where the repeated game rejects.
\item \emph{Common marginals and diagonal weight}: the game's two marginals agree,
  with common marginal $\nu$, and the almost-synchronous transport requires
  $\mu(x,x)\geq\kappa\nu(x)$ for every $x$, with the same
  $\kappa > 0$ along the whole recursion. This is the hypothesis of
  Theorem~\ref{thm:almost-sync}, without which the two values of a game are unrelated
  (Remark~\ref{rem:almost-sync-hypothesis}).
\end{enumerate}
The second property is emphatically not one of repeated games: the $k$-fold repetition of
a game with diagonal weight $\kappa$ has diagonal weight $\kappa^k$. A synchronous route
would therefore have to establish it for the constructions of introspection,
oracularization and answer reduction --- never inheriting it from the game being
compressed, and never applying the transport to a repeated game --- and in~\cite{JNVWY20}
both properties come from the typed structure of the samplers, which this blueprint does
not carry (Section~\ref{sec:departures}). Neither is established here, and neither needs
to be.
\end{remark}

\begin{definition}[$\lambda$-bounded verifier in the ambient model]
\label{def:lambda-bounded}
\uses{def:normal-verifier}
\lean{MIPRE.Verifier.IsBounded}
\leanok
For $\lambda\in\N$, an ambient normal form verifier $\verifier=(\sampler,\decider)$
is \emph{$\lambda$-bounded} if $\abs{\verifier}\leq\lambda$ and, for every $n\geq2$,
$s(n)\leq n^\lambda$ and both programs halt within cost
$n^\lambda(\abs{d}+1)^\lambda$ on every encoded input $(n,d)$. Equivalently, both
\texttt{Sampler.TimeBoundAt} and \texttt{Decider.TimeBoundAt} hold at parameters
$(n,n^\lambda,\lambda)$. The parameter $\lambda$ controls
the coefficient's growth in $n$ and bounds the polynomial degree in the auxiliary input
size. The dimension bound is a separate clause; neither it nor an absolute-time bound
is inferred from the relative cost predicates.
\end{definition}

\begin{remark}[Comparison with source boundedness]
\label{rem:source-lambda-bounded}
\uses{def:lambda-bounded,rem:source-normal-verifier}
The paper's \texttt{def:lambda} instead requires
$\TIME_\sampler(n),\TIME_\decider(n)\leq n^\lambda$ in the absolute-time model,
for $n\geq2$, and $\abs{\verifier}\leq\lambda$~\cite{AuditSync26}. The source dimension
bound follows from writing the sampler's output and from the decider's initial checks.
This condition and the marked ambient definition above are different contracts; their
comparison requires the normalization obligation below.
\end{remark}

\begin{remark}[The runtime-contract bridge remains an obligation]
\label{rem:runtime-contract-bridge}
\uses{def:sampler,def:decider,def:normal-verifier,def:lambda-bounded}
The paper's transformation contracts use absolute simulation budgets and, where stated,
standard timeout form. They cannot be applied to \texttt{Verifier.IsBounded} solely by
reinterpreting $\TIME$: even on inputs of size $O(N^\lambda)$ the relative budget
$N^\lambda(\abs{d}+1)^\lambda$ can be much larger than $N^\lambda$. A normalization
bridge must supply bounded-prefix behavior and the required timeout form, preserve the
sampler distribution and the chosen finite game, and transfer the relevant completeness
witnesses. It must keep the original answer cut by rejecting longer answers rather than
silently enlarging the alphabet when it increases the simulation budget. Any new absolute
bound, for example $N^{p(\lambda)}$, must then be propagated through the transformation
parameters. This is a pending bridge between the source and ambient interfaces; the
existing conditional Lean results do not construct it.
\end{remark}

\begin{definition}[The computable class $\MIPstar$ used in Lean]
\label{def:mipstar}
\uses{def:q-value,def:game-description}
\lean{MIPRE.MIPStar}
\leanok
A language $L\subseteq\{0,1\}^*$ is in the class denoted $\MIPstar$ here if there exists
a computable map $g$ from bit strings to \texttt{HaltingGameValue.GameData}
(Definition~\ref{def:game-description}) such that for every string $z$,
\begin{equation*}
  z\in L\ \Longrightarrow\ \valstar(\game_{g(z)})=1,
  \qquad
  z\notin L\ \Longrightarrow\ \valstar(\game_{g(z)})\leq\tfrac12.
\end{equation*}
The game $\game_{g(z)}$ is the explicit game \texttt{(g z).game}. Computability is the
Mathlib \texttt{Computable} predicate; there is no polynomial-time bound on $g$ in this
definition. This is exactly \texttt{MIPRE.MIPStar}, the class used by the checked
inclusion in $\RE$ and the conditional main theorem.
\end{definition}

\begin{remark}[Comparison with the source complexity class]
\label{rem:source-mipstar}
\ledgernode{1.1.7, 1.1.7.1}
\uses{def:mipstar,rem:source-normal-verifier}
The paper's $\MIP^*_{1,1/2}(2,1)$ uses a sampler and decider running in time polynomial
in the input length, with the same perfect-completeness and $\frac12$-soundness gap
(\texttt{games.tex}~\cite{AuditSync26}). That polynomial-time verifier convention is not
part of the marked class above. Explicitly tabulating a finite game is a computable
operation but may take much more than polynomial time. Moreover, general bipartite games
must be represented in \texttt{GameData} with value-preserving transport; merely forcing
unequal answers to equal questions to lose can destroy a perfect strategy. The checked
halting reduction supplies its particular explicit descriptions, while the source
polynomial-time complexity-class claim requires its own efficient verifier construction.
\end{remark}
